%%C040001 p

本章包含以下案例：


%%C040002 li

\begin{itemize}
\item 条件编译源代码


%%C040003 li

\item 用间接层实现预处理器的字符串化与连接


%%C040004 li

\item 用static\_assert做编译期断言检查


%%C040005 li

\item 用enable\_if条件编译类与函数


%%C040006 li

\item 用constexpr if在编译期选择分支


%%C040007 li

\item 用属性向编译器提供元数据

\end{itemize}


%%C040008 h2

\section{引言}


%%C040009 p

在C++中，编译是把源代码转换成机器码、组织进目标文件，再把这些目标文件链接起来生成可执行文件的过程。编译器实际上一次只处理一个文件，该文件由预处理器根据单个源文件及其包含的所有头文件生成。不过，这只是对编译代码时所发生事情的一种过度简化。本章讨论与预处理和编译相关的主题，重点介绍进行条件编译的各种方法，同时也涉及其他现代主题，例如用属性提供由实现定义的语言扩展。


%%C040010 h2

\section{条件编译源代码}


%%C040011 p

条件编译是一种简单的机制，它让开发者能够维护单一的代码库，但只把代码的某些部分纳入编译，以生成不同的可执行文件，通常是为了在不同的平台、硬件上运行，或者依赖不同的库或库版本。常见的例子包括根据编译器、平台（x86、x64、ARM等）、配置（调试版或发布版）或任何用户定义的具体条件来决定使用还是忽略某段代码。本案例将介绍条件编译的工作方式。


%%C040012 h3

\subsection{准备工作}


%%C040013 p

条件编译是一项因多种目的而被广泛使用的技术。本案例将看几个示例并解释它们的工作原理。这项技术绝不限于这些示例。就本案例的范围而言，我们只考虑三大编译器：GCC、Clang和VC++。


%%C040014 h3

\subsection{具体做法}


%%C040015 p

要条件编译代码的某些部分，请使用\texttt{\#if}、\texttt{\#ifdef}和\texttt{\#ifndef}指令（配合\texttt{\#elif}、\texttt{\#else}和\texttt{\#endif}指令）。条件编译的一般形式如下：


%%C040016 code

\begin{codeblockbg}
\begin{Verbatim}[commandchars=\\\{\},breaklines=true,breakanywhere=true,fontsize=\codefont,xleftmargin=17bp]
\PYG{c+cp}{\PYGZsh{}}\PYG{c+cp}{if condition1}
\PYG{+w}{  }\PYG{n}{text}
\PYG{c+cp}{\PYGZsh{}}\PYG{c+cp}{elif condition2}
\PYG{+w}{  }\PYG{n}{text}
\PYG{c+cp}{\PYGZsh{}}\PYG{c+cp}{elif condition3}
\PYG{+w}{  }\PYG{n}{text}
\PYG{c+cp}{\PYGZsh{}}\PYG{c+cp}{else}
\PYG{+w}{  }\PYG{n}{text}
\PYG{c+cp}{\PYGZsh{}}\PYG{c+cp}{endif}
\end{Verbatim}
\end{codeblockbg}


%%C040017 p

要为条件编译定义宏，可以使用以下两种方式之一：


%%C040018 li

\begin{itemize}
\item 在源代码中使用\texttt{\#define}指令：


%%C040019 code

\begin{codeblockbg}
\begin{Verbatim}[commandchars=\\\{\},breaklines=true,breakanywhere=true,fontsize=\codefont,xleftmargin=17bp]
\PYG{c+cp}{\PYGZsh{}}\PYG{c+cp}{define DEBUG\PYGZus{}PRINTS}
\PYG{c+cp}{\PYGZsh{}}\PYG{c+cp}{define VERBOSITY\PYGZus{}LEVEL 5}
\end{Verbatim}
\end{codeblockbg}


%%C040020 li

\item 各编译器特有的命令行选项。最常用编译器的示例如下：


%%C040021 li

\begin{itemize}
\item 对于Visual C++，使用\texttt{/\allowbreak{}Dname}或\texttt{/\allowbreak{}Dname=value}（其中\texttt{/\allowbreak{}Dname}等价于\texttt{/\allowbreak{}Dname=1}），

\end{itemize}

\end{itemize}


%%C040022 p

例如，\myttbr{cl /\allowbreak{}DVERBOSITY\_\allowbreak{}LEVEL=5}。


%%C040023 li

\begin{itemize}
\item 对于GCC和Clang，使用\texttt{-\allowbreak{}D name}或\texttt{-\allowbreak{}D name=value}（其中\texttt{-\allowbreak{}D} \texttt{name}等价于\texttt{-\allowbreak{}D name=1}），

\end{itemize}


%%C040024 p

例如，\myttbr{gcc -\allowbreak{}D VERBOSITY\_\allowbreak{}LEVEL=5}。


%%C040025 p

以下是条件编译的典型示例：


%%C040026 li

\begin{itemize}
\item 用头文件保护符避免重复定义：


%%C040027 code

\begin{codeblockbg}
\begin{Verbatim}[commandchars=\\\{\},breaklines=true,breakanywhere=true,fontsize=\codefont,xleftmargin=17bp]
\PYG{c+cp}{\PYGZsh{}}\PYG{c+cp}{if !defined(\PYGZus{}UNIQUE\PYGZus{}NAME\PYGZus{})}
\PYG{c+cp}{\PYGZsh{}}\PYG{c+cp}{define \PYGZus{}UNIQUE\PYGZus{}NAME\PYGZus{}}
\PYG{k}{class}\PYG{+w}{ }\PYG{n+nc}{foo}\PYG{+w}{ }\PYG{p}{\PYGZob{}}\PYG{+w}{ }\PYG{p}{\PYGZcb{}}\PYG{p}{;}
\PYG{c+cp}{\PYGZsh{}}\PYG{c+cp}{endif}
\end{Verbatim}
\end{codeblockbg}


%%C040028 li

\item 跨平台应用中的编译器相关代码。以下示例把编译器的名称打印到控制台：


%%C040029 code

\begin{codeblockbg}
\begin{Verbatim}[commandchars=\\\{\},breaklines=true,breakanywhere=true,fontsize=\codefont,xleftmargin=17bp]
\PYG{k+kt}{void}\PYG{+w}{ }\PYG{n+nf}{show\PYGZus{}compiler}\PYG{p}{(}\PYG{p}{)}
\PYG{p}{\PYGZob{}}
\PYG{+w}{  }\PYG{c+cp}{\PYGZsh{}}\PYG{c+cp}{if defined \PYGZus{}MSC\PYGZus{}VER}
\PYG{+w}{    }\PYG{n}{std}\PYG{o}{:}\PYG{o}{:}\PYG{n}{cout}\PYG{+w}{ }\PYG{o}{\PYGZlt{}}\PYG{o}{\PYGZlt{}}\PYG{+w}{ }\PYG{l+s}{\PYGZdq{}}\PYG{l+s}{Visual C++}\PYG{l+s}{\PYGZdq{}}\PYG{+w}{ }\PYG{o}{\PYGZlt{}}\PYG{o}{\PYGZlt{}}\PYG{+w}{ }\PYG{n}{std}\PYG{o}{:}\PYG{o}{:}\PYG{n}{endl}\PYG{p}{;}
\PYG{+w}{  }\PYG{c+cp}{\PYGZsh{}}\PYG{c+cp}{elif defined \PYGZus{}\PYGZus{}clang\PYGZus{}\PYGZus{}}
\PYG{+w}{    }\PYG{n}{std}\PYG{o}{:}\PYG{o}{:}\PYG{n}{cout}\PYG{+w}{ }\PYG{o}{\PYGZlt{}}\PYG{o}{\PYGZlt{}}\PYG{+w}{ }\PYG{l+s}{\PYGZdq{}}\PYG{l+s}{Clang}\PYG{l+s}{\PYGZdq{}}\PYG{+w}{ }\PYG{o}{\PYGZlt{}}\PYG{o}{\PYGZlt{}}\PYG{+w}{ }\PYG{n}{std}\PYG{o}{:}\PYG{o}{:}\PYG{n}{endl}\PYG{p}{;}
\PYG{+w}{  }\PYG{c+cp}{\PYGZsh{}}\PYG{c+cp}{elif defined \PYGZus{}\PYGZus{}GNUG\PYGZus{}\PYGZus{}}
\PYG{+w}{    }\PYG{n}{std}\PYG{o}{:}\PYG{o}{:}\PYG{n}{cout}\PYG{+w}{ }\PYG{o}{\PYGZlt{}}\PYG{o}{\PYGZlt{}}\PYG{+w}{ }\PYG{l+s}{\PYGZdq{}}\PYG{l+s}{GCC}\PYG{l+s}{\PYGZdq{}}\PYG{+w}{ }\PYG{o}{\PYGZlt{}}\PYG{o}{\PYGZlt{}}\PYG{+w}{ }\PYG{n}{std}\PYG{o}{:}\PYG{o}{:}\PYG{n}{endl}\PYG{p}{;}
\PYG{+w}{  }\PYG{c+cp}{\PYGZsh{}}\PYG{c+cp}{else}
\PYG{+w}{    }\PYG{n}{std}\PYG{o}{:}\PYG{o}{:}\PYG{n}{cout}\PYG{+w}{ }\PYG{o}{\PYGZlt{}}\PYG{o}{\PYGZlt{}}\PYG{+w}{ }\PYG{l+s}{\PYGZdq{}}\PYG{l+s}{Unknown compiler}\PYG{l+s}{\PYGZdq{}}\PYG{+w}{ }\PYG{o}{\PYGZlt{}}\PYG{o}{\PYGZlt{}}\PYG{+w}{ }\PYG{n}{std}\PYG{o}{:}\PYG{o}{:}\PYG{n}{endl}\PYG{p}{;}
\PYG{+w}{  }\PYG{c+cp}{\PYGZsh{}}\PYG{c+cp}{endif}
\PYG{p}{\PYGZcb{}}
\end{Verbatim}
\end{codeblockbg}


%%C040030 li

\item 面向多种架构的目标平台相关代码，例如为多个编译器和架构条件编译代码：


%%C040031 code

\begin{codeblockbg}
\begin{Verbatim}[commandchars=\\\{\},breaklines=true,breakanywhere=true,fontsize=\codefont,xleftmargin=17bp]
\PYG{k+kt}{void}\PYG{+w}{ }\PYG{n+nf}{show\PYGZus{}architecture}\PYG{p}{(}\PYG{p}{)}
\PYG{p}{\PYGZob{}}
\PYG{+w}{  }\PYG{c+cp}{\PYGZsh{}}\PYG{c+cp}{if defined \PYGZus{}MSC\PYGZus{}VER}
\PYG{n}{std}\PYG{o}{:}\PYG{o}{:}\PYG{n}{cout}\PYG{+w}{ }\PYG{o}{\PYGZlt{}}\PYG{o}{\PYGZlt{}}
\PYG{c+cp}{\PYGZsh{}}\PYG{c+cp}{if defined \PYGZus{}M\PYGZus{}X64}
\PYG{+w}{    }\PYG{l+s}{\PYGZdq{}}\PYG{l+s}{AMD64}\PYG{l+s}{\PYGZdq{}}
\PYG{+w}{  }\PYG{c+cp}{\PYGZsh{}}\PYG{c+cp}{elif defined \PYGZus{}M\PYGZus{}IX86}
\PYG{+w}{    }\PYG{l+s}{\PYGZdq{}}\PYG{l+s}{INTEL x86}\PYG{l+s}{\PYGZdq{}}
\PYG{+w}{  }\PYG{c+cp}{\PYGZsh{}}\PYG{c+cp}{elif defined \PYGZus{}M\PYGZus{}ARM}
\PYG{+w}{    }\PYG{l+s}{\PYGZdq{}}\PYG{l+s}{ARM}\PYG{l+s}{\PYGZdq{}}
\PYG{+w}{  }\PYG{c+cp}{\PYGZsh{}}\PYG{c+cp}{else}
\PYG{+w}{    }\PYG{l+s}{\PYGZdq{}}\PYG{l+s}{unknown}\PYG{l+s}{\PYGZdq{}}
\PYG{+w}{  }\PYG{c+cp}{\PYGZsh{}}\PYG{c+cp}{endif}
\PYG{o}{\PYGZlt{}}\PYG{o}{\PYGZlt{}}\PYG{+w}{ }\PYG{n}{std}\PYG{o}{:}\PYG{o}{:}\PYG{n}{endl}\PYG{p}{;}
\PYG{c+cp}{\PYGZsh{}}\PYG{c+cp}{elif defined \PYGZus{}\PYGZus{}clang\PYGZus{}\PYGZus{} || \PYGZus{}\PYGZus{}GNUG\PYGZus{}\PYGZus{}}
\PYG{n}{std}\PYG{o}{:}\PYG{o}{:}\PYG{n}{cout}\PYG{+w}{ }\PYG{o}{\PYGZlt{}}\PYG{o}{\PYGZlt{}}
\PYG{+w}{  }\PYG{c+cp}{\PYGZsh{}}\PYG{c+cp}{if defined \PYGZus{}\PYGZus{}amd64\PYGZus{}\PYGZus{}}
\PYG{+w}{    }\PYG{l+s}{\PYGZdq{}}\PYG{l+s}{AMD64}\PYG{l+s}{\PYGZdq{}}
\PYG{+w}{  }\PYG{c+cp}{\PYGZsh{}}\PYG{c+cp}{elif defined \PYGZus{}\PYGZus{}i386\PYGZus{}\PYGZus{}}
\PYG{+w}{    }\PYG{l+s}{\PYGZdq{}}\PYG{l+s}{INTEL x86}\PYG{l+s}{\PYGZdq{}}
\PYG{+w}{  }\PYG{c+cp}{\PYGZsh{}}\PYG{c+cp}{elif defined \PYGZus{}\PYGZus{}arm\PYGZus{}\PYGZus{}}
\PYG{+w}{    }\PYG{l+s}{\PYGZdq{}}\PYG{l+s}{ARM}\PYG{l+s}{\PYGZdq{}}
\PYG{+w}{  }\PYG{c+cp}{\PYGZsh{}}\PYG{c+cp}{else}
\PYG{+w}{    }\PYG{l+s}{\PYGZdq{}}\PYG{l+s}{unknown}\PYG{l+s}{\PYGZdq{}}
\PYG{+w}{  }\PYG{c+cp}{\PYGZsh{}}\PYG{c+cp}{endif}
\PYG{o}{\PYGZlt{}}\PYG{o}{\PYGZlt{}}\PYG{+w}{ }\PYG{n}{std}\PYG{o}{:}\PYG{o}{:}\PYG{n}{endl}\PYG{p}{;}
\PYG{+w}{  }\PYG{c+cp}{\PYGZsh{}}\PYG{c+cp}{else}
\PYG{+w}{  }\PYG{c+cp}{\PYGZsh{}}\PYG{c+cp}{error Unknown compiler}
\PYG{+w}{  }\PYG{c+cp}{\PYGZsh{}}\PYG{c+cp}{endif}
\PYG{p}{\PYGZcb{}}
\end{Verbatim}
\end{codeblockbg}


%%C040032 li

\item 配置相关的代码，例如为调试版和发布版条件编译代码：


%%C040033 code

\begin{codeblockbg}
\begin{Verbatim}[commandchars=\\\{\},breaklines=true,breakanywhere=true,fontsize=\codefont,xleftmargin=17bp]
\PYG{k+kt}{void}\PYG{+w}{ }\PYG{n+nf}{show\PYGZus{}configuration}\PYG{p}{(}\PYG{p}{)}
\PYG{p}{\PYGZob{}}
\PYG{+w}{  }\PYG{n}{std}\PYG{o}{:}\PYG{o}{:}\PYG{n}{cout}\PYG{+w}{ }\PYG{o}{\PYGZlt{}}\PYG{o}{\PYGZlt{}}
\PYG{+w}{  }\PYG{c+cp}{\PYGZsh{}}\PYG{c+cp}{ifdef \PYGZus{}DEBUG}
\PYG{+w}{    }\PYG{l+s}{\PYGZdq{}}\PYG{l+s}{debug}\PYG{l+s}{\PYGZdq{}}
\PYG{+w}{  }\PYG{c+cp}{\PYGZsh{}}\PYG{c+cp}{else}
\PYG{+w}{    }\PYG{l+s}{\PYGZdq{}}\PYG{l+s}{release}\PYG{l+s}{\PYGZdq{}}
\PYG{+w}{  }\PYG{c+cp}{\PYGZsh{}}\PYG{c+cp}{endif}
\PYG{+w}{  }\PYG{o}{\PYGZlt{}}\PYG{o}{\PYGZlt{}}\PYG{+w}{ }\PYG{n}{std}\PYG{o}{:}\PYG{o}{:}\PYG{n}{endl}\PYG{p}{;}
\PYG{p}{\PYGZcb{}}
\end{Verbatim}
\end{codeblockbg}

\end{itemize}


%%C040034 h3

\subsection{工作原理}


%%C040035 p

当你使用预处理器指令\texttt{\#if}、\texttt{\#ifndef}、\texttt{\#ifdef}、\texttt{\#elif}、\texttt{\#else}和\texttt{\#endif}时，编译器最多只会选择一个分支，把该分支的主体纳入翻译单元参与编译。这些指令的主体可以是任意文本，包括其他预处理器指令。适用以下规则：


%%C040036 li

\begin{itemize}
\item \texttt{\#if}、\texttt{\#ifdef}和\texttt{\#ifndef}必须由\texttt{\#endif}配对。


%%C040037 li

\item \texttt{\#if}指令可以有多个\texttt{\#elif}指令，但只能有一个\texttt{\#else}，而且它必须是\texttt{\#endif}之前的最后一个。


%%C040038 li

\item \texttt{\#if}、\texttt{\#ifdef}、\texttt{\#ifndef}、\texttt{\#elif}、\texttt{\#else}和\texttt{\#endif}可以嵌套。


%%C040039 li

\item \texttt{\#if}指令要求常量表达式，而\texttt{\#ifdef}和\texttt{\#ifndef}要求标识符。


%%C040040 li

\item 运算符\texttt{defined}可用于预处理器常量表达式，但只能出现在\texttt{\#if}和\texttt{\#elif}指令中。


%%C040041 li

\item 如果\texttt{identifier}已定义，则\myttbr{defined(\allowbreak{}identifier)}被视为\texttt{true}，否则被视为\texttt{false}。


%%C040042 li

\item 定义为空文本的标识符也被视为已定义。


%%C040043 li

\item \texttt{\#ifdef identifier}等价于\myttbr{\#if defined(\allowbreak{}identifier)}。


%%C040044 li

\item \myttbr{\#ifndef identifier}等价于\myttbr{\#if !defined(\allowbreak{}identifier)}。


%%C040045 li

\item \myttbr{defined(\allowbreak{}identifier)}与\texttt{defined identifier}等价。

\end{itemize}


%%C040046 p

头文件保护符是最常见的条件编译形式之一。这项技术用来防止头文件的内容被多次编译（不过每次仍会扫描头文件，以确定应当包含哪些内容）。由于头文件经常被多个源文件包含，如果每个包含它的翻译单元都对其进行编译，就会为同一批符号产生多个定义，这是错误的。因此，头文件中的代码会按上一节示例所示的方式防止重复编译。以给出的示例来说，其工作方式是：如果宏\texttt{\_\allowbreak{}UNIQUE\_\allowbreak{}NAME\_}未定义，那么\texttt{\#if}指令之后直到\texttt{\#endif}的代码就会被纳入翻译单元并编译。此时会用\texttt{\#define}指令定义宏\texttt{\_\allowbreak{}UNIQUE\_\allowbreak{}NAME\_}。下次该头文件被包含进翻译单元时，宏\texttt{\_\allowbreak{}UNIQUE\_\allowbreak{}NAME\_}已经定义，\texttt{\#if}指令主体中的代码就不会被纳入翻译单元，因此不会再被编译。


%%C040047 note

\begin{myNote}
要注意，宏的名称必须在整个应用程序中唯一，否则只会编译第一个使用该宏的头文件中的代码，其他使用同一名称的头文件中的代码将被忽略。
\end{myNote}


%%C040048 p

条件编译的另一个重要例子是跨平台代码，它需要兼顾不同的编译器和架构，通常是Intel x86、AMD64或ARM之一。不过，编译器会为可能的平台定义自己的宏。\textit{具体做法}一节的示例展示了如何为多个编译器和架构条件编译代码。


%%C040049 note

\begin{myNote}
要注意，在上述示例中我们只考虑了几种架构。实践中，识别同一种架构可能有多个宏。在代码中使用这类宏之前，务必阅读各编译器的文档。
\end{myNote}


%%C040050 p

配置相关的代码同样通过宏和条件编译来处理。GCC和Clang等编译器不会为调试配置（使用\texttt{-\allowbreak{}g}选项时）定义任何特殊宏。Visual C++确实会为调试配置定义\texttt{\_\allowbreak{}DEBUG}，这一点已在\textit{具体做法}一节的最后一个示例中展示。对于其他编译器，你必须显式定义一个宏来标识这样的调试配置。


%%C040051 h3

\subsection{延伸阅读}


%%C040052 li

\begin{itemize}
\item \textit{用间接层实现预处理器的字符串化与连接}

\end{itemize}


%%C040053 h2

\section{用间接层实现预处理器的字符串化与连接}


%%C040054 p

C++预处理器提供两个运算符，分别用于把标识符转换成字符串，以及把标识符连接起来。第一个是\texttt{operator \#}，称为\textit{字符串化}运算符；第二个是\texttt{operator \#\#}，称为\textit{记号粘贴}、\textit{合并}或\textit{连接}运算符。虽然它们的使用仅限于某些特定情形，但理解其工作方式仍然很重要。


%%C040055 h3

\subsection{准备工作}


%%C040056 p

本案例要求你会用预处理器指令\texttt{\#define}定义宏。


%%C040057 h3

\subsection{具体做法}


%%C040058 p

要使用预处理器的\texttt{operator\#}从标识符创建字符串，请按以下模式操作：


%%C040059 p

定义一个接受一个参数的辅助宏，它展开为\texttt{\#}后跟该


%%C040060 p

参数：


%%C040061 code

\begin{codeblockbg}
\begin{Verbatim}[commandchars=\\\{\},breaklines=true,breakanywhere=true,fontsize=\codefont,xleftmargin=17bp]
\PYG{c+cp}{\PYGZsh{}}\PYG{c+cp}{define MAKE\PYGZus{}STR2(x) \PYGZsh{}x}
\end{Verbatim}
\end{codeblockbg}


%%C040062 p

定义你要使用的宏，它接受一个参数，展开为该


%%C040063 p

辅助宏：


%%C040064 code

\begin{codeblockbg}
\begin{Verbatim}[commandchars=\\\{\},breaklines=true,breakanywhere=true,fontsize=\codefont,xleftmargin=17bp]
\PYG{c+cp}{\PYGZsh{}}\PYG{c+cp}{define MAKE\PYGZus{}STR(x) MAKE\PYGZus{}STR2(x)}
\end{Verbatim}
\end{codeblockbg}


%%C040065 p

要使用预处理器的\texttt{operator\#\#}把标识符连接起来，请按以下模式操作：


%%C040066 p

定义一个带有一个或多个参数的辅助宏，这些参数使用记号粘贴


%%C040067 p

运算符\texttt{\#\#}来连接参数：


%%C040068 code

\begin{codeblockbg}
\begin{Verbatim}[commandchars=\\\{\},breaklines=true,breakanywhere=true,fontsize=\codefont,xleftmargin=17bp]
\PYG{c+cp}{\PYGZsh{}}\PYG{c+cp}{define MERGE2(x, y)    x\PYGZsh{}\PYGZsh{}y}
\end{Verbatim}
\end{codeblockbg}


%%C040069 p

通过这个辅助宏来定义你要使用的宏：


%%C040070 code

\begin{codeblockbg}
\begin{Verbatim}[commandchars=\\\{\},breaklines=true,breakanywhere=true,fontsize=\codefont,xleftmargin=17bp]
\PYG{c+cp}{\PYGZsh{}}\PYG{c+cp}{define MERGE(x, y)     MERGE2(x, y)}
\end{Verbatim}
\end{codeblockbg}


%%C040071 h3

\subsection{工作原理}


%%C040072 p

为了理解它们的工作方式，来看看前面定义的\texttt{MAKE\_\allowbreak{}STR}和\texttt{MAKE\_\allowbreak{}STR2}宏。把它们用于任意文本时，都会生成包含该文本的字符串。下例展示了如何用这两个宏定义包含文本\textit{“sample”}的字符串。


%%C040073 code

\begin{codeblockbg}
\begin{Verbatim}[commandchars=\\\{\},breaklines=true,breakanywhere=true,fontsize=\codefont,xleftmargin=17bp]
\PYG{n}{std}\PYG{o}{:}\PYG{o}{:}\PYG{n}{string}\PYG{+w}{ }\PYG{n}{s1}\PYG{+w}{ }\PYG{p}{\PYGZob{}}\PYG{+w}{ }\PYG{n}{MAKE\PYGZus{}STR}\PYG{p}{(}\PYG{n}{sample}\PYG{p}{)}\PYG{+w}{ }\PYG{p}{\PYGZcb{}}\PYG{p}{;}\PYG{+w}{  }\PYG{c+c1}{// s1 = \PYGZdq{}sample\PYGZdq{}}
\PYG{n}{std}\PYG{o}{:}\PYG{o}{:}\PYG{n}{string}\PYG{+w}{ }\PYG{n}{s2}\PYG{+w}{ }\PYG{p}{\PYGZob{}}\PYG{+w}{ }\PYG{n}{MAKE\PYGZus{}STR2}\PYG{p}{(}\PYG{n}{sample}\PYG{p}{)}\PYG{+w}{ }\PYG{p}{\PYGZcb{}}\PYG{p}{;}\PYG{+w}{ }\PYG{c+c1}{// s2 = \PYGZdq{}sample\PYGZdq{}}
\end{Verbatim}
\end{codeblockbg}


%%C040074 p

另一方面，把宏作为参数传入时，结果就不同了。在下例中，\texttt{NUMBER}是展开为整数\texttt{42}的宏。把它作为\texttt{MAKE\_\allowbreak{}STR}的参数时，确实生成了字符串\textit{“42”}；但作为\texttt{MAKE\_\allowbreak{}STR2}的参数时，生成的字符串是\textit{“NUMBER”}：


%%C040075 code

\begin{codeblockbg}
\begin{Verbatim}[commandchars=\\\{\},breaklines=true,breakanywhere=true,fontsize=\codefont,xleftmargin=17bp]
\PYG{c+cp}{\PYGZsh{}}\PYG{c+cp}{define NUMBER 42}
\PYG{n}{std}\PYG{o}{:}\PYG{o}{:}\PYG{n}{string}\PYG{+w}{ }\PYG{n}{s3}\PYG{+w}{ }\PYG{p}{\PYGZob{}}\PYG{+w}{ }\PYG{n}{MAKE\PYGZus{}STR}\PYG{p}{(}\PYG{n}{NUMBER}\PYG{p}{)}\PYG{+w}{ }\PYG{p}{\PYGZcb{}}\PYG{p}{;}\PYG{+w}{    }\PYG{c+c1}{// s3 = \PYGZdq{}42\PYGZdq{}}
\PYG{n}{std}\PYG{o}{:}\PYG{o}{:}\PYG{n}{string}\PYG{+w}{ }\PYG{n}{s4}\PYG{+w}{ }\PYG{p}{\PYGZob{}}\PYG{+w}{ }\PYG{n}{MAKE\PYGZus{}STR2}\PYG{p}{(}\PYG{n}{NUMBER}\PYG{p}{)}\PYG{+w}{ }\PYG{p}{\PYGZcb{}}\PYG{p}{;}\PYG{+w}{   }\PYG{c+c1}{// s4 = \PYGZdq{}NUMBER\PYGZdq{}}
\end{Verbatim}
\end{codeblockbg}


%%C040076 p

C++标准对类函数宏中的实参替换规定了以下规则（第16.3.1节）：


%%C040077 p

\textit{在识别出类函数宏调用所用的实参之后，} \textit{就会进行实参替换。替换列表中的形参，除非前面带有} \textit{\#或\#\#预处理记号，或者后面跟着\#\#预处理记号（见下文），否则} \textit{在其中的所有宏都展开之后，该形参会被替换为对应的实参。} \textit{在被替换之前，每个实参的预处理记号都会被完全} \textit{宏替换，就好像它们构成了预处理文件的其余部分；没有其他预处理} \textit{记号可用。}


%%C040078 p

这段话的意思是：宏实参在替换进宏体之前会先展开，但宏体中的形参前面或后面是运算符\texttt{\#}或\texttt{\#\#}的情况除外。因此，会出现以下情况：


%%C040079 li

\begin{itemize}
\item 对于\texttt{MAKE\_\allowbreak{}STR2(\allowbreak{}NUMBER)}，替换列表中的\texttt{NUMBER}形参前面是\texttt{\#}，因此把实参替换进宏体之前它不会展开；于是替换后得到\texttt{\#NUMBER}，最终变成\textit{“NUMBER”}。


%%C040080 li

\item 对于\texttt{MAKE\_\allowbreak{}STR(\allowbreak{}NUMBER)}，替换列表是\texttt{MAKE\_\allowbreak{}STR2(\allowbreak{}NUMBER)}，其中没有\texttt{\#}或\texttt{\#\#}；因此\texttt{NUMBER}形参在被替换之前就换成了对应的实参\texttt{42}。结果是\texttt{MAKE\_\allowbreak{}STR2(\allowbreak{}42)}，它会被再次扫描，展开后变成\textit{“42”}。

\end{itemize}


%%C040081 p

同样的处理规则也适用于使用记号粘贴运算符的宏。因此，为了确保字符串化宏和连接宏在所有情况下都能正常工作，请始终采用本案例所述的间接层模式。


%%C040082 p

记号粘贴运算符通常用于把重复代码提取出来的宏中，以避免一遍又一遍地显式书写相同内容。下面这个简单的示例展示了记号粘贴运算符的实际用途；给定一组类，我们希望提供为每个类创建实例的工厂方法：


%%C040083 code

\begin{codeblockbg}
\begin{Verbatim}[commandchars=\\\{\},breaklines=true,breakanywhere=true,fontsize=\codefont,xleftmargin=17bp]
\PYG{c+cp}{\PYGZsh{}}\PYG{c+cp}{define DECL\PYGZus{}MAKE(x)    DECL\PYGZus{}MAKE2(x)}
\PYG{c+cp}{\PYGZsh{}}\PYG{c+cp}{define DECL\PYGZus{}MAKE2(x)   x* make\PYGZsh{}\PYGZsh{}\PYGZus{}\PYGZsh{}\PYGZsh{}x() \PYGZob{} return new x(); \PYGZcb{}}
\PYG{k}{struct}\PYG{+w}{ }\PYG{n+nc}{bar}\PYG{+w}{ }\PYG{p}{\PYGZob{}}\PYG{p}{\PYGZcb{}}\PYG{p}{;}
\PYG{k}{struct}\PYG{+w}{ }\PYG{n+nc}{foo}\PYG{+w}{ }\PYG{p}{\PYGZob{}}\PYG{p}{\PYGZcb{}}\PYG{p}{;}
\PYG{n}{DECL\PYGZus{}MAKE}\PYG{p}{(}\PYG{n}{foo}\PYG{p}{)}
\PYG{n}{DECL\PYGZus{}MAKE}\PYG{p}{(}\PYG{n}{bar}\PYG{p}{)}
\PYG{k}{auto}\PYG{+w}{ }\PYG{n}{f}\PYG{+w}{ }\PYG{o}{=}\PYG{+w}{ }\PYG{n}{make\PYGZus{}foo}\PYG{p}{(}\PYG{p}{)}\PYG{p}{;}\PYG{+w}{ }\PYG{c+c1}{// f 是 foo*}
\PYG{k}{auto}\PYG{+w}{ }\PYG{n}{b}\PYG{+w}{ }\PYG{o}{=}\PYG{+w}{ }\PYG{n}{make\PYGZus{}bar}\PYG{p}{(}\PYG{p}{)}\PYG{p}{;}\PYG{+w}{ }\PYG{c+c1}{// b 是 bar*}
\end{Verbatim}
\end{codeblockbg}


%%C040084 p

熟悉Windows平台的人可能用过\texttt{\_\allowbreak{}T}（或\texttt{\_\allowbreak{}TEXT}）宏来声明字符串字面量，它们会被翻译成Unicode或ANSI字符串（涵盖单字节和多字节字符字符串）：


%%C040085 code

\begin{codeblockbg}
\begin{Verbatim}[commandchars=\\\{\},breaklines=true,breakanywhere=true,fontsize=\codefont,xleftmargin=17bp]
\PYG{k}{auto}\PYG{+w}{ }\PYG{n}{text}\PYG{p}{\PYGZob{}}\PYG{+w}{ }\PYG{n}{\PYGZus{}T}\PYG{p}{(}\PYG{l+s}{\PYGZdq{}}\PYG{l+s}{sample}\PYG{l+s}{\PYGZdq{}}\PYG{p}{)}\PYG{+w}{ }\PYG{p}{\PYGZcb{}}\PYG{p}{;}\PYG{+w}{ }\PYG{c+c1}{// text 要么是 \PYGZdq{}sample\PYGZdq{}，要么是 L\PYGZdq{}sample\PYGZdq{}}
\end{Verbatim}
\end{codeblockbg}


%%C040086 p

Windows SDK对\texttt{\_\allowbreak{}T}宏的定义如下。要注意，当定义了\texttt{\_\allowbreak{}UNICODE}时，该宏会用记号粘贴运算符把\texttt{L}前缀与传给宏的实际字符串连接起来：


%%C040087 code

\begin{codeblockbg}
\begin{Verbatim}[commandchars=\\\{\},breaklines=true,breakanywhere=true,fontsize=\codefont,xleftmargin=17bp]
\PYG{c+cp}{\PYGZsh{}}\PYG{c+cp}{ifdef \PYGZus{}UNICODE}
\PYG{c+cp}{\PYGZsh{}}\PYG{c+cp}{define \PYGZus{}\PYGZus{}T(x)   L \PYGZsh{}\PYGZsh{} x}
\PYG{c+cp}{\PYGZsh{}}\PYG{c+cp}{else}
\PYG{c+cp}{\PYGZsh{}}\PYG{c+cp}{define \PYGZus{}\PYGZus{}T(x)   x}
\PYG{c+cp}{\PYGZsh{}}\PYG{c+cp}{endif}
\PYG{c+cp}{\PYGZsh{}}\PYG{c+cp}{define \PYGZus{}T(x)    \PYGZus{}\PYGZus{}T(x)}
\PYG{c+cp}{\PYGZsh{}}\PYG{c+cp}{define \PYGZus{}TEXT(x) \PYGZus{}\PYGZus{}T(x)}
\end{Verbatim}
\end{codeblockbg}


%%C040088 p

乍看之下，让一个宏去调用另一个宏似乎没有必要，但正是这一层间接，才使得\texttt{\#}和\texttt{\#\#}运算符能与其他宏配合工作，正如本案例中所见。


%%C040089 h3

\subsection{延伸阅读}


%%C040090 li

\begin{itemize}
\item \textit{条件编译源代码}

\end{itemize}


%%C040091 h2

\section{用static\_assert做编译期断言检查}


%%C040092 p

在C++中，既可以做运行期断言检查，也可以做编译期断言检查，以确保代码中的特定条件成立。运行期断言的缺点是验证得晚：只有在程序运行、并且控制流到达该处时才会检查。当条件依赖运行期数据时，别无他法；但如果不是这种情况，就应当优先采用编译期断言检查。使用编译期断言，编译器能够在开发阶段早期就以错误的形式通知你某个条件不满足。不过，这类断言只有在条件能在编译期求值时才能使用。在C++11中，编译期断言通过\texttt{static\_\allowbreak{}assert}完成。


%%C040093 h3

\subsection{准备工作}


%%C040094 p

静态断言检查最常见的用途是模板元编程，用它来验证模板类型的前置条件是否满足（例如类型是否为POD类型、是否可复制构造、是否为引用类型等）。另一个典型例子是确保类型（或对象）具有预期的大小。


%%C040095 h3

\subsection{具体做法}


%%C040096 p

使用\texttt{static\_\allowbreak{}assert}声明来确保不同作用域中的条件成立：


%%C040097 li

\begin{itemize}
\item \textit{命名空间}：本示例中，我们验证类\texttt{item}的大小始终为16：


%%C040098 code

\begin{codeblockbg}
\begin{Verbatim}[commandchars=\\\{\},breaklines=true,breakanywhere=true,fontsize=\codefont,xleftmargin=17bp]
\PYG{k}{struct}\PYG{+w}{ }\PYG{n+nc}{alignas}\PYG{p}{(}\PYG{l+m+mi}{8}\PYG{p}{)}\PYG{+w}{ }\PYG{n}{item}
\PYG{p}{\PYGZob{}}
\PYG{+w}{  }\PYG{k+kt}{int}\PYG{+w}{      }\PYG{n}{id}\PYG{p}{;}
\PYG{+w}{  }\PYG{k+kt}{bool}\PYG{+w}{     }\PYG{n}{active}\PYG{p}{;}
\PYG{+w}{  }\PYG{k+kt}{double}\PYG{+w}{   }\PYG{n}{value}\PYG{p}{;}
\PYG{p}{\PYGZcb{}}\PYG{p}{;}
\PYG{k}{static\PYGZus{}assert}\PYG{p}{(}\PYG{k}{sizeof}\PYG{p}{(}\PYG{n}{item}\PYG{p}{)}\PYG{+w}{ }\PYG{o}{=}\PYG{o}{=}\PYG{+w}{ }\PYG{l+m+mi}{16}\PYG{p}{,}
\PYG{+w}{              }\PYG{l+s}{\PYGZdq{}}\PYG{l+s}{size of item must be 16 bytes}\PYG{l+s}{\PYGZdq{}}\PYG{p}{)}\PYG{p}{;}
\end{Verbatim}
\end{codeblockbg}


%%C040099 li

\item \textit{类}：本示例中，我们验证\texttt{pod\_\allowbreak{}wrapper}只能用于POD类型：


%%C040100 code

\begin{codeblockbg}
\begin{Verbatim}[commandchars=\\\{\},breaklines=true,breakanywhere=true,fontsize=\codefont,xleftmargin=17bp]
\PYG{k}{template}\PYG{+w}{ }\PYG{o}{\PYGZlt{}}\PYG{k}{typename}\PYG{+w}{ }\PYG{n+nc}{T}\PYG{o}{\PYGZgt{}}
\PYG{k}{class}\PYG{+w}{ }\PYG{n+nc}{pod\PYGZus{}wrapper}
\PYG{p}{\PYGZob{}}
\PYG{+w}{  }\PYG{k}{static\PYGZus{}assert}\PYG{p}{(}\PYG{n}{std}\PYG{o}{:}\PYG{o}{:}\PYG{n}{is\PYGZus{}pod}\PYG{o}{\PYGZlt{}}\PYG{n}{T}\PYG{o}{\PYGZgt{}}\PYG{o}{:}\PYG{o}{:}\PYG{n}{value}\PYG{p}{,}
\PYG{+w}{                }\PYG{l+s}{\PYGZdq{}}\PYG{l+s}{POD type expected!}\PYG{l+s}{\PYGZdq{}}\PYG{p}{)}\PYG{p}{;}
\PYG{+w}{  }\PYG{n}{T}\PYG{+w}{ }\PYG{n}{value}\PYG{p}{;}
\PYG{p}{\PYGZcb{}}\PYG{p}{;}
\PYG{k}{struct}\PYG{+w}{ }\PYG{n+nc}{point}
\PYG{p}{\PYGZob{}}
\PYG{+w}{  }\PYG{k+kt}{int}\PYG{+w}{ }\PYG{n}{x}\PYG{p}{;}
\PYG{+w}{  }\PYG{k+kt}{int}\PYG{+w}{ }\PYG{n}{y}\PYG{p}{;}
\PYG{p}{\PYGZcb{}}\PYG{p}{;}
\PYG{n}{pod\PYGZus{}wrapper}\PYG{o}{\PYGZlt{}}\PYG{k+kt}{int}\PYG{o}{\PYGZgt{}}\PYG{+w}{         }\PYG{n}{w1}\PYG{p}{;}\PYG{+w}{ }\PYG{c+c1}{// 正确}
\PYG{n}{pod\PYGZus{}wrapper}\PYG{o}{\PYGZlt{}}\PYG{n}{point}\PYG{o}{\PYGZgt{}}\PYG{+w}{       }\PYG{n}{w2}\PYG{p}{;}\PYG{+w}{ }\PYG{c+c1}{// 正确}
\PYG{n}{pod\PYGZus{}wrapper}\PYG{o}{\PYGZlt{}}\PYG{n}{std}\PYG{o}{:}\PYG{o}{:}\PYG{n}{string}\PYG{o}{\PYGZgt{}}\PYG{+w}{ }\PYG{n}{w3}\PYG{p}{;}\PYG{+w}{ }\PYG{c+c1}{// 错误：需要POD类型}
\end{Verbatim}
\end{codeblockbg}


%%C040101 li

\item \textit{块（函数）}：本示例中，我们验证函数模板只能接受整型实参：


%%C040102 code

\begin{codeblockbg}
\begin{Verbatim}[commandchars=\\\{\},breaklines=true,breakanywhere=true,fontsize=\codefont,xleftmargin=17bp]
\PYG{k}{template}\PYG{o}{\PYGZlt{}}\PYG{k}{typename}\PYG{+w}{ }\PYG{n+nc}{T}\PYG{o}{\PYGZgt{}}\PYG{+w}{ }\PYG{k}{auto}
\PYG{n}{mul}\PYG{p}{(}\PYG{n}{T}\PYG{+w}{ }\PYG{k}{const}\PYG{+w}{ }\PYG{n}{a}\PYG{p}{,}\PYG{+w}{ }\PYG{n}{T}\PYG{+w}{ }\PYG{k}{const}\PYG{+w}{ }\PYG{n}{b}\PYG{p}{)}
\PYG{p}{\PYGZob{}}
\PYG{+w}{  }\PYG{k}{static\PYGZus{}assert}\PYG{p}{(}\PYG{n}{std}\PYG{o}{:}\PYG{o}{:}\PYG{n}{is\PYGZus{}integral}\PYG{o}{\PYGZlt{}}\PYG{n}{T}\PYG{o}{\PYGZgt{}}\PYG{o}{:}\PYG{o}{:}\PYG{n}{value}\PYG{p}{,}
\PYG{+w}{                }\PYG{l+s}{\PYGZdq{}}\PYG{l+s}{Integral type expected}\PYG{l+s}{\PYGZdq{}}\PYG{p}{)}\PYG{p}{;}
\PYG{+w}{  }\PYG{k}{return}\PYG{+w}{ }\PYG{n}{a}\PYG{+w}{ }\PYG{o}{*}\PYG{+w}{ }\PYG{n}{b}\PYG{p}{;}
\PYG{p}{\PYGZcb{}}
\PYG{k}{auto}\PYG{+w}{ }\PYG{n}{v1}\PYG{+w}{ }\PYG{o}{=}\PYG{+w}{ }\PYG{n}{mul}\PYG{p}{(}\PYG{l+m+mi}{1}\PYG{p}{,}\PYG{+w}{ }\PYG{l+m+mi}{2}\PYG{p}{)}\PYG{p}{;}\PYG{+w}{       }\PYG{c+c1}{// 正确}
\PYG{k}{auto}\PYG{+w}{ }\PYG{n}{v2}\PYG{+w}{ }\PYG{o}{=}\PYG{+w}{ }\PYG{n}{mul}\PYG{p}{(}\PYG{l+m+mf}{12.0}\PYG{p}{,}\PYG{+w}{ }\PYG{l+m+mf}{42.5}\PYG{p}{)}\PYG{p}{;}\PYG{+w}{ }\PYG{c+c1}{// 错误：需要整型}
\end{Verbatim}
\end{codeblockbg}

\end{itemize}


%%C040103 h3

\subsection{工作原理}


%%C040104 p

\texttt{static\_\allowbreak{}assert}本质上是一条声明，但它不引入新的名字。这类声明的形式如下：


%%C040105 code

\begin{codeblockbg}
\begin{Verbatim}[commandchars=\\\{\},breaklines=true,breakanywhere=true,fontsize=\codefont,xleftmargin=17bp]
\PYG{k}{static\PYGZus{}assert}\PYG{p}{(}\PYG{n}{condition}\PYG{p}{,}\PYG{+w}{ }\PYG{n}{message}\PYG{p}{)}\PYG{p}{;}
\end{Verbatim}
\end{codeblockbg}


%%C040106 p

条件必须能在编译期转换为布尔值，消息必须是字符串字面量。从C++17起，消息是可选的。

%%C040107 p

当\texttt{static\_\allowbreak{}assert}声明中的条件求值为\texttt{true}时，什么也不会发生。当条件求值为\texttt{false}时，编译器会产生一个错误，并在其中包含所指定的消息（如果给出了的话）。


%%C040108 h3

\subsection{延伸阅读}


%%C040109 li

\begin{itemize}
\item \textit{用enable\_if条件编译类与函数}


%%C040110 li

\item \textit{用constexpr if在编译期选择分支}

\end{itemize}


%%C040111 h2

\section{用enable\_if条件编译类与函数}


%%C040112 p

模板元编程是C++的一项强大特性，它让我们能够编写适用于任何类型的泛型类和泛型函数。但有时这恰恰是个问题，因为语言并没有定义任何机制，用来对可以替换模板参数的类型施加约束。不过，我们仍然可以通过元编程技巧，并借助一条称为\textit{替换失败} \textit{并非错误}的规则（简称\textit{SFINAE}）来做到这一点。本案例将着重讲解如何为模板实现类型约束。


%%C040113 h3

\subsection{准备工作}


%%C040114 p

多年来，开发者一直把通常称为\texttt{enable\_\allowbreak{}if}的类模板与SFINAE结合使用，以实现对模板类型的约束。\texttt{enable\_\allowbreak{}if}这一族模板已成为C++11标准的一部分，其实现如下：


%%C040115 code

\begin{codeblockbg}
\begin{Verbatim}[commandchars=\\\{\},breaklines=true,breakanywhere=true,fontsize=\codefont,xleftmargin=17bp]
\PYG{k}{template}\PYG{o}{\PYGZlt{}}\PYG{k+kt}{bool}\PYG{+w}{ }\PYG{n}{Test}\PYG{p}{,}\PYG{+w}{ }\PYG{k}{class}\PYG{+w}{ }\PYG{n+nc}{T}\PYG{+w}{ }\PYG{o}{=}\PYG{+w}{ }\PYG{k+kt}{void}\PYG{o}{\PYGZgt{}}
\PYG{k}{struct}\PYG{+w}{ }\PYG{n+nc}{enable\PYGZus{}if}
\PYG{p}{\PYGZob{}}\PYG{p}{\PYGZcb{}}\PYG{p}{;}
\PYG{k}{template}\PYG{o}{\PYGZlt{}}\PYG{k}{class}\PYG{+w}{ }\PYG{n+nc}{T}\PYG{o}{\PYGZgt{}}
\PYG{k}{struct}\PYG{+w}{ }\PYG{n+nc}{enable\PYGZus{}if}\PYG{o}{\PYGZlt{}}\PYG{n+nb}{true}\PYG{p}{,}\PYG{+w}{ }\PYG{n}{T}\PYG{o}{\PYGZgt{}}
\PYG{p}{\PYGZob{}}
\PYG{+w}{  }\PYG{k}{typedef}\PYG{+w}{ }\PYG{n}{T}\PYG{+w}{ }\PYG{n}{type}\PYG{p}{;}
\PYG{p}{\PYGZcb{}}\PYG{p}{;}
\end{Verbatim}
\end{codeblockbg}


%%C040116 p

要使用\texttt{std::\allowbreak{}enable\_\allowbreak{}if}，必须包含头文件\texttt{<\allowbreak{}type\_\allowbreak{}traits>}。


%%C040117 h3

\subsection{具体做法}


%%C040118 p

\texttt{std::\allowbreak{}enable\_\allowbreak{}if}可以用在多种作用域中，以达到不同的目的；考虑下面这些例子：


%%C040119 li

\begin{itemize}
\item 用在类模板参数上，使类模板只对满足指定条件的类型启用：


%%C040120 code

\begin{codeblockbg}
\begin{Verbatim}[commandchars=\\\{\},breaklines=true,breakanywhere=true,fontsize=\codefont,xleftmargin=17bp]
\PYG{k}{template}\PYG{+w}{ }\PYG{o}{\PYGZlt{}}\PYG{k}{typename}\PYG{+w}{ }\PYG{n+nc}{T}\PYG{p}{,}
\PYG{+w}{          }\PYG{k}{typename}\PYG{+w}{ }\PYG{o}{=}\PYG{+w}{ }\PYG{k}{typename}
\PYG{+w}{             }\PYG{n+nc}{std}\PYG{o}{:}\PYG{o}{:}\PYG{n}{enable\PYGZus{}if}\PYG{o}{\PYGZlt{}}\PYG{n}{std}\PYG{o}{:}\PYG{o}{:}\PYG{n}{is\PYGZus{}pod}\PYG{o}{\PYGZlt{}}\PYG{n}{T}\PYG{o}{\PYGZgt{}}\PYG{o}{:}\PYG{o}{:}\PYG{n}{value}\PYG{p}{,}\PYG{+w}{ }\PYG{n}{T}\PYG{o}{\PYGZgt{}}\PYG{o}{:}\PYG{o}{:}\PYG{n}{type}\PYG{o}{\PYGZgt{}}
\PYG{k}{class}\PYG{+w}{ }\PYG{n+nc}{pod\PYGZus{}wrapper}
\PYG{p}{\PYGZob{}}
\PYG{+w}{  }\PYG{n}{T}\PYG{+w}{ }\PYG{n}{value}\PYG{p}{;}
\PYG{p}{\PYGZcb{}}\PYG{p}{;}
\PYG{k}{struct}\PYG{+w}{ }\PYG{n+nc}{point}
\PYG{p}{\PYGZob{}}
\PYG{+w}{  }\PYG{k+kt}{int}\PYG{+w}{ }\PYG{n}{x}\PYG{p}{;}
\PYG{+w}{  }\PYG{k+kt}{int}\PYG{+w}{ }\PYG{n}{y}\PYG{p}{;}
\PYG{p}{\PYGZcb{}}\PYG{p}{;}
\PYG{n}{pod\PYGZus{}wrapper}\PYG{o}{\PYGZlt{}}\PYG{k+kt}{int}\PYG{o}{\PYGZgt{}}\PYG{+w}{         }\PYG{n}{w1}\PYG{p}{;}\PYG{+w}{ }\PYG{c+c1}{// OK}
\PYG{n}{pod\PYGZus{}wrapper}\PYG{o}{\PYGZlt{}}\PYG{n}{point}\PYG{o}{\PYGZgt{}}\PYG{+w}{       }\PYG{n}{w2}\PYG{p}{;}\PYG{+w}{ }\PYG{c+c1}{// OK}
\PYG{n}{pod\PYGZus{}wrapper}\PYG{o}{\PYGZlt{}}\PYG{n}{std}\PYG{o}{:}\PYG{o}{:}\PYG{n}{string}\PYG{o}{\PYGZgt{}}\PYG{+w}{ }\PYG{n}{w3}\PYG{p}{;}\PYG{+w}{ }\PYG{c+c1}{// 错误：模板参数太少}
\end{Verbatim}
\end{codeblockbg}


%%C040121 li

\item 用在函数模板参数、函数参数或函数返回类型上，使函数模板只对满足指定条件的类型启用：


%%C040122 code

\begin{codeblockbg}
\begin{Verbatim}[commandchars=\\\{\},breaklines=true,breakanywhere=true,fontsize=\codefont,xleftmargin=17bp]
\PYG{k}{template}\PYG{o}{\PYGZlt{}}\PYG{k}{typename}\PYG{+w}{ }\PYG{n+nc}{T}\PYG{p}{,}
\PYG{+w}{         }\PYG{k}{typename}\PYG{+w}{ }\PYG{o}{=}\PYG{+w}{ }\PYG{k}{typename}\PYG{+w}{ }\PYG{n+nc}{std}\PYG{o}{:}\PYG{o}{:}\PYG{n}{enable\PYGZus{}if}\PYG{o}{\PYGZlt{}}
\PYG{+w}{            }\PYG{n}{std}\PYG{o}{:}\PYG{o}{:}\PYG{n}{is\PYGZus{}integral}\PYG{o}{\PYGZlt{}}\PYG{n}{T}\PYG{o}{\PYGZgt{}}\PYG{o}{:}\PYG{o}{:}\PYG{n}{value}\PYG{p}{,}\PYG{+w}{ }\PYG{n}{T}\PYG{o}{\PYGZgt{}}\PYG{o}{:}\PYG{o}{:}\PYG{n}{type}\PYG{o}{\PYGZgt{}}
\PYG{k}{auto}\PYG{+w}{ }\PYG{n}{mul}\PYG{p}{(}\PYG{n}{T}\PYG{+w}{ }\PYG{k}{const}\PYG{+w}{ }\PYG{n}{a}\PYG{p}{,}\PYG{+w}{ }\PYG{n}{T}\PYG{+w}{ }\PYG{k}{const}\PYG{+w}{ }\PYG{n}{b}\PYG{p}{)}
\PYG{p}{\PYGZob{}}
\PYG{+w}{  }\PYG{k}{return}\PYG{+w}{ }\PYG{n}{a}\PYG{+w}{ }\PYG{o}{*}\PYG{+w}{ }\PYG{n}{b}\PYG{p}{;}
\PYG{p}{\PYGZcb{}}
\PYG{k}{auto}\PYG{+w}{ }\PYG{n}{v1}\PYG{+w}{ }\PYG{o}{=}\PYG{+w}{ }\PYG{n}{mul}\PYG{p}{(}\PYG{l+m+mi}{1}\PYG{p}{,}\PYG{+w}{ }\PYG{l+m+mi}{2}\PYG{p}{)}\PYG{p}{;}\PYG{+w}{     }\PYG{c+c1}{// OK}
\PYG{k}{auto}\PYG{+w}{ }\PYG{n}{v2}\PYG{+w}{ }\PYG{o}{=}\PYG{+w}{ }\PYG{n}{mul}\PYG{p}{(}\PYG{l+m+mf}{1.0}\PYG{p}{,}\PYG{+w}{ }\PYG{l+m+mf}{2.0}\PYG{p}{)}\PYG{p}{;}
\PYG{c+c1}{// 错误：找不到匹配的重载函数}
\end{Verbatim}
\end{codeblockbg}

\end{itemize}


%%C040123 p

为了简化使用\texttt{std::\allowbreak{}enable\_\allowbreak{}if}时写出的那堆繁杂代码，我们可以借助别名模板，定义两个别名\texttt{EnableIf}和\texttt{DisabledIf}：


%%C040124 code

\begin{codeblockbg}
\begin{Verbatim}[commandchars=\\\{\},breaklines=true,breakanywhere=true,fontsize=\codefont,xleftmargin=17bp]
\PYG{k}{template}\PYG{+w}{ }\PYG{o}{\PYGZlt{}}\PYG{k}{typename}\PYG{+w}{ }\PYG{n+nc}{Test}\PYG{p}{,}\PYG{+w}{ }\PYG{k}{typename}\PYG{+w}{ }\PYG{n+nc}{T}\PYG{+w}{ }\PYG{o}{=}\PYG{+w}{ }\PYG{k+kt}{void}\PYG{o}{\PYGZgt{}}
\PYG{k}{using}\PYG{+w}{ }\PYG{n}{EnableIf}\PYG{+w}{ }\PYG{o}{=}\PYG{+w}{ }\PYG{k}{typename}\PYG{+w}{ }\PYG{n+nc}{std}\PYG{o}{:}\PYG{o}{:}\PYG{n}{enable\PYGZus{}if}\PYG{o}{\PYGZlt{}}\PYG{n}{Test}\PYG{o}{:}\PYG{o}{:}\PYG{n}{value}\PYG{p}{,}\PYG{+w}{ }\PYG{n}{T}\PYG{o}{\PYGZgt{}}\PYG{o}{:}\PYG{o}{:}\PYG{n}{type}\PYG{p}{;}
\PYG{k}{template}\PYG{+w}{ }\PYG{o}{\PYGZlt{}}\PYG{k}{typename}\PYG{+w}{ }\PYG{n+nc}{Test}\PYG{p}{,}\PYG{+w}{ }\PYG{k}{typename}\PYG{+w}{ }\PYG{n+nc}{T}\PYG{+w}{ }\PYG{o}{=}\PYG{+w}{ }\PYG{k+kt}{void}\PYG{o}{\PYGZgt{}}
\PYG{k}{using}\PYG{+w}{ }\PYG{n}{DisableIf}\PYG{+w}{ }\PYG{o}{=}\PYG{+w}{ }\PYG{k}{typename}\PYG{+w}{ }\PYG{n+nc}{std}\PYG{o}{:}\PYG{o}{:}\PYG{n}{enable\PYGZus{}if}\PYG{o}{\PYGZlt{}}\PYG{o}{!}\PYG{n}{Test}\PYG{o}{:}\PYG{o}{:}\PYG{n}{value}\PYG{p}{,}\PYG{+w}{ }\PYG{n}{T}\PYG{o}{\PYGZgt{}}\PYG{o}{:}\PYG{o}{:}\PYG{n}{type}\PYG{p}{;}
\end{Verbatim}
\end{codeblockbg}


%%C040125 p

基于这两个别名模板，下面的定义与前面展示的定义等价：


%%C040126 code

\begin{codeblockbg}
\begin{Verbatim}[commandchars=\\\{\},breaklines=true,breakanywhere=true,fontsize=\codefont,xleftmargin=17bp]
\PYG{k}{template}\PYG{+w}{ }\PYG{o}{\PYGZlt{}}\PYG{k}{typename}\PYG{+w}{ }\PYG{n+nc}{T}\PYG{p}{,}\PYG{+w}{ }\PYG{k}{typename}\PYG{+w}{ }\PYG{o}{=}\PYG{+w}{ }\PYG{n}{EnableIf}\PYG{o}{\PYGZlt{}}\PYG{n}{std}\PYG{o}{:}\PYG{o}{:}\PYG{n}{is\PYGZus{}pod}\PYG{o}{\PYGZlt{}}\PYG{n}{T}\PYG{o}{\PYGZgt{}}\PYG{o}{\PYGZgt{}}\PYG{o}{\PYGZgt{}}
\PYG{k}{class}\PYG{+w}{ }\PYG{n+nc}{pod\PYGZus{}wrapper}
\PYG{p}{\PYGZob{}}
\PYG{+w}{  }\PYG{n}{T}\PYG{+w}{ }\PYG{n}{value}\PYG{p}{;}
\PYG{p}{\PYGZcb{}}\PYG{p}{;}
\PYG{k}{template}\PYG{o}{\PYGZlt{}}\PYG{k}{typename}\PYG{+w}{ }\PYG{n+nc}{T}\PYG{p}{,}\PYG{+w}{ }\PYG{k}{typename}\PYG{+w}{ }\PYG{o}{=}\PYG{+w}{ }\PYG{n}{EnableIf}\PYG{o}{\PYGZlt{}}\PYG{n}{std}\PYG{o}{:}\PYG{o}{:}\PYG{n}{is\PYGZus{}integral}\PYG{o}{\PYGZlt{}}\PYG{n}{T}\PYG{o}{\PYGZgt{}}\PYG{o}{\PYGZgt{}}\PYG{o}{\PYGZgt{}}
\PYG{k}{auto}\PYG{+w}{ }\PYG{n}{mul}\PYG{p}{(}\PYG{n}{T}\PYG{+w}{ }\PYG{k}{const}\PYG{+w}{ }\PYG{n}{a}\PYG{p}{,}\PYG{+w}{ }\PYG{n}{T}\PYG{+w}{ }\PYG{k}{const}\PYG{+w}{ }\PYG{n}{b}\PYG{p}{)}
\PYG{p}{\PYGZob{}}
\PYG{+w}{  }\PYG{k}{return}\PYG{+w}{ }\PYG{n}{a}\PYG{+w}{ }\PYG{o}{*}\PYG{+w}{ }\PYG{n}{b}\PYG{p}{;}
\PYG{p}{\PYGZcb{}}
\end{Verbatim}
\end{codeblockbg}


%%C040127 h3

\subsection{工作原理}


%%C040128 p

\texttt{std::\allowbreak{}enable\_\allowbreak{}if}之所以能工作，是因为编译器在执行重载决议时应用了SFINAE规则。在解释\texttt{std::\allowbreak{}enable\_\allowbreak{}if}如何工作之前，我们应当先快速看看SFINAE到底是什么。


%%C040129 p

当编译器遇到函数调用时，它需要构建一组可能的重载，并根据该次函数调用的实参从中选出最佳匹配。在构建这个重载集时，编译器也会对函数模板求值，并且必须把指定或推导出的类型替换到模板实参中。根据SFINAE，当替换失败时，编译器不应当报错，而应当把这个函数模板从重载集中移除，然后继续。


%%C040130 note

\begin{myNote}[tip]
标准规定了一组同样属于SFINAE错误的类型错误和表达式错误。其中包括：试图创建\texttt{void}数组或长度为零的数组，试图创建对\texttt{void}的引用，试图创建参数类型为\texttt{void}的函数类型，以及试图在模板实参表达式中或在函数声明所用的表达式中执行无效转换。完整的例外清单请查阅C++标准或其他资料。
\end{myNote}


%%C040131 p

考虑名为\texttt{func()}的函数的下面两个重载。第一个重载是函数模板，它有一个类型为\texttt{T::\allowbreak{}data\_\allowbreak{}type}的参数；这意味着它只能用带有名为\texttt{data\_\allowbreak{}type}的内部类型的类型来实例化。第二个重载是函数，它有一个类型为\texttt{int}的参数：


%%C040132 code

\begin{codeblockbg}
\begin{Verbatim}[commandchars=\\\{\},breaklines=true,breakanywhere=true,fontsize=\codefont,xleftmargin=17bp]
\PYG{k}{template}\PYG{+w}{ }\PYG{o}{\PYGZlt{}}\PYG{k}{typename}\PYG{+w}{ }\PYG{n+nc}{T}\PYG{o}{\PYGZgt{}}
\PYG{k+kt}{void}\PYG{+w}{ }\PYG{n}{func}\PYG{p}{(}\PYG{k}{typename}\PYG{+w}{ }\PYG{n+nc}{T}\PYG{o}{:}\PYG{o}{:}\PYG{n}{data\PYGZus{}type}\PYG{+w}{ }\PYG{k}{const}\PYG{+w}{ }\PYG{n}{a}\PYG{p}{)}
\PYG{p}{\PYGZob{}}\PYG{+w}{ }\PYG{n}{std}\PYG{o}{:}\PYG{o}{:}\PYG{n}{cout}\PYG{+w}{ }\PYG{o}{\PYGZlt{}}\PYG{o}{\PYGZlt{}}\PYG{+w}{ }\PYG{l+s}{\PYGZdq{}}\PYG{l+s}{func\PYGZlt{}\PYGZgt{}}\PYG{l+s}{\PYGZdq{}}\PYG{+w}{ }\PYG{o}{\PYGZlt{}}\PYG{o}{\PYGZlt{}}\PYG{+w}{ }\PYG{n}{std}\PYG{o}{:}\PYG{o}{:}\PYG{n}{endl}\PYG{p}{;}\PYG{+w}{ }\PYG{p}{\PYGZcb{}}
\PYG{k+kt}{void}\PYG{+w}{ }\PYG{n}{func}\PYG{p}{(}\PYG{k+kt}{int}\PYG{+w}{ }\PYG{k}{const}\PYG{+w}{ }\PYG{n}{a}\PYG{p}{)}
\PYG{p}{\PYGZob{}}\PYG{+w}{ }\PYG{n}{std}\PYG{o}{:}\PYG{o}{:}\PYG{n}{cout}\PYG{+w}{ }\PYG{o}{\PYGZlt{}}\PYG{o}{\PYGZlt{}}\PYG{+w}{ }\PYG{l+s}{\PYGZdq{}}\PYG{l+s}{func}\PYG{l+s}{\PYGZdq{}}\PYG{+w}{ }\PYG{o}{\PYGZlt{}}\PYG{o}{\PYGZlt{}}\PYG{+w}{ }\PYG{n}{std}\PYG{o}{:}\PYG{o}{:}\PYG{n}{endl}\PYG{p}{;}\PYG{+w}{ }\PYG{p}{\PYGZcb{}}
\PYG{k}{template}\PYG{+w}{ }\PYG{o}{\PYGZlt{}}\PYG{k}{typename}\PYG{+w}{ }\PYG{n+nc}{T}\PYG{o}{\PYGZgt{}}
\PYG{k}{struct}\PYG{+w}{ }\PYG{n+nc}{some\PYGZus{}type}
\PYG{p}{\PYGZob{}}
\PYG{+w}{  }\PYG{k}{using}\PYG{+w}{ }\PYG{n}{data\PYGZus{}type}\PYG{+w}{ }\PYG{o}{=}\PYG{+w}{ }\PYG{n}{T}\PYG{p}{;}
\PYG{p}{\PYGZcb{}}\PYG{p}{;}
\end{Verbatim}
\end{codeblockbg}


%%C040133 p

如果编译器遇到诸如\texttt{func(\allowbreak{}42)}这样的调用，它就必须找到一个能接受\texttt{int}实参的重载。当它构建重载集、并用所提供的模板实参替换模板参数时，得到的结果\myttbr{void func(\allowbreak{}int::\allowbreak{}data\_\allowbreak{}type} \texttt{const)}是无效的，因为\texttt{int}并没有\texttt{data\_\allowbreak{}type}成员。由于SFINAE，编译器不会报错并停止，而只是忽略这个重载并继续。随后它找到\myttbr{void func(\allowbreak{}int const)}，这就是它将调用的最佳（也是唯一）匹配。


%%C040134 p

如果编译器遇到诸如\myttbr{func<\allowbreak{}some\_\allowbreak{}type<\allowbreak{}int>>(\allowbreak{}42)}这样的调用，那么它会构建一个包含\myttbr{void func(\allowbreak{}some\_\allowbreak{}type<\allowbreak{}int>::\allowbreak{}data\_\allowbreak{}type const>}和\texttt{void} \texttt{func(\allowbreak{}int const)}的重载集，此例中的最佳匹配是第一个重载；这一次不涉及SFINAE。


%%C040135 p

另一方面，如果编译器遇到诸如\texttt{func(\allowbreak{}“string”s)}这样的调用，它同样依靠SFINAE忽略该函数模板，因为\texttt{std::\allowbreak{}basic\_\allowbreak{}string}也没有\texttt{value\_\allowbreak{}type}成员。不过这一次重载集中没有任何与字符串实参匹配的候选；因此，程序非良构，编译器报错并停止。


%%C040136 p

类模板\myttbr{enable\_\allowbreak{}if<\allowbreak{}bool, T>}没有任何成员，但它的偏特化\myttbr{enable\_\allowbreak{}if<\allowbreak{}true, T>}确实有一个名为\texttt{type}的内部类型，它是\texttt{T}的同义词。当作为第一个实参传给\texttt{enable\_\allowbreak{}if}的编译期表达式求值为\texttt{true}时，内部成员\texttt{type}就可用，否则不可用。


%%C040137 p

对于\textit{具体做法}一节中\texttt{mul()}函数的最后一个定义，当编译器遇到诸如\texttt{mul(\allowbreak{}1, 2)}这样的调用时，它会尝试用\texttt{int}替换模板参数\texttt{T}；由于\texttt{int}是整数类型，\myttbr{std::\allowbreak{}is\_\allowbreak{}integral<\allowbreak{}T>}求值为\texttt{true}，因此会实例化一个定义了名为\texttt{type}的内部类型的\texttt{enable\_\allowbreak{}if}特化。结果，别名模板\texttt{EnableIf}成为该类型的同义词，而这个类型是\texttt{void}（来自表达式\texttt{typename T = void}）。最终得到的是函数模板\myttbr{int mul<\allowbreak{}int, void>(\allowbreak{}int a, int b)}，它可以用所提供的实参调用。


%%C040138 p

另一方面，当编译器遇到诸如\texttt{mul(\allowbreak{}1.\allowbreak{}0, 2.\allowbreak{}0)}这样的调用时，它会尝试用\texttt{double}替换模板参数\texttt{T}。然而这并不是整数类型；因此\texttt{std::\allowbreak{}enable\_\allowbreak{}if}中的条件求值为\texttt{false}，类模板也就不定义内部成员\texttt{type}。这会导致替换错误，但根据SFINAE，编译器不会报错，而是继续处理。然而由于找不到其他重载，最终没有任何\texttt{mul()}函数可供调用。于是程序被视为非良构，编译器报错停止。


%%C040139 p

类模板\texttt{pod\_\allowbreak{}wrapper}会遇到类似的情况。它有两个模板类型参数：第一个是要包装的实际POD类型，第二个是\texttt{enable\_\allowbreak{}if}与\texttt{is\_\allowbreak{}pod}替换的结果。如果该类型是POD类型（如\texttt{pod\_\allowbreak{}wrapper<\allowbreak{}int>}），那么来自\texttt{enable\_\allowbreak{}if}的内部成员\texttt{type}就存在，并由它来替换第二个模板类型参数。然而如果内部成员\texttt{type}不是POD类型（如\myttbr{pod\_\allowbreak{}wrapper<\allowbreak{}std::\allowbreak{}string>}），那么内部成员\texttt{type}就没有定义，替换失败，从而产生诸如\textit{“模板参数太少”}这样的错误。


%%C040140 h3

\subsection{更多内容}


%%C040141 p

\texttt{static\_\allowbreak{}assert}和\texttt{std::\allowbreak{}enable\_\allowbreak{}if}都可以用来达到相同的目的。事实上，在上一个案例\textit{用static\_assert做编译期断言检查}中，我们就定义了同样的类模板\texttt{pod\_\allowbreak{}wrapper}和函数模板\texttt{mul()}。对于这些例子，\texttt{static\_\allowbreak{}assert}似乎是更好的方案，因为编译器会给出更好的错误消息（前提是你在\texttt{static\_\allowbreak{}assert}声明中写了相关的消息）。不过两者的工作方式差别很大，并不是用来互相替代的。


%%C040142 p

\texttt{static\_\allowbreak{}assert}不依赖SFINAE，而是在重载决议完成之后才起作用。断言失败的结果就是编译错误。相反，\texttt{std::\allowbreak{}enable\_\allowbreak{}if}用于把候选从重载集中移除，并不会触发编译错误（前提是不出现标准为SFINAE规定的那些例外）。SFINAE之后真正可能出现的错误，是重载集为空，使得某个特定的函数调用无法完成，从而令程序非良构。


%%C040143 p

为了理解\texttt{static\_\allowbreak{}assert}与配合SFINAE的\texttt{std::\allowbreak{}enable\_\allowbreak{}if}之间的区别，假设我们想要两个函数重载：一个用于整数类型的实参，另一个用于整数类型以外的任何实参。使用\texttt{static\_\allowbreak{}assert}，可以这样写（注意第二个重载上那个占位的第二类型参数是必需的，用来定义两个不同的重载，否则我们就只是同一个函数的两个定义）：


%%C040144 code

\begin{codeblockbg}
\begin{Verbatim}[commandchars=\\\{\},breaklines=true,breakanywhere=true,fontsize=\codefont,xleftmargin=17bp]
\PYG{k}{template}\PYG{+w}{ }\PYG{o}{\PYGZlt{}}\PYG{k}{typename}\PYG{+w}{ }\PYG{n+nc}{T}\PYG{o}{\PYGZgt{}}
\PYG{k}{auto}\PYG{+w}{ }\PYG{n}{compute}\PYG{p}{(}\PYG{n}{T}\PYG{+w}{ }\PYG{k}{const}\PYG{+w}{ }\PYG{n}{a}\PYG{p}{,}\PYG{+w}{ }\PYG{n}{T}\PYG{+w}{ }\PYG{k}{const}\PYG{+w}{ }\PYG{n}{b}\PYG{p}{)}
\PYG{p}{\PYGZob{}}
\PYG{+w}{  }\PYG{k}{static\PYGZus{}assert}\PYG{p}{(}\PYG{n}{std}\PYG{o}{:}\PYG{o}{:}\PYG{n}{is\PYGZus{}integral}\PYG{o}{\PYGZlt{}}\PYG{n}{T}\PYG{o}{\PYGZgt{}}\PYG{p}{,}
\PYG{+w}{                }\PYG{l+s}{\PYGZdq{}}\PYG{l+s}{An integral type expected}\PYG{l+s}{\PYGZdq{}}\PYG{p}{)}\PYG{p}{;}
\PYG{+w}{  }\PYG{k}{return}\PYG{+w}{ }\PYG{n}{a}\PYG{+w}{ }\PYG{o}{+}\PYG{+w}{ }\PYG{n}{b}\PYG{p}{;}
\PYG{p}{\PYGZcb{}}
\PYG{k}{template}\PYG{+w}{ }\PYG{o}{\PYGZlt{}}\PYG{k}{typename}\PYG{+w}{ }\PYG{n+nc}{T}\PYG{p}{,}\PYG{+w}{ }\PYG{k}{typename}\PYG{+w}{ }\PYG{o}{=}\PYG{+w}{ }\PYG{k+kt}{void}\PYG{o}{\PYGZgt{}}
\PYG{k}{auto}\PYG{+w}{ }\PYG{n}{compute}\PYG{p}{(}\PYG{n}{T}\PYG{+w}{ }\PYG{k}{const}\PYG{+w}{ }\PYG{n}{a}\PYG{p}{,}\PYG{+w}{ }\PYG{n}{T}\PYG{+w}{ }\PYG{k}{const}\PYG{+w}{ }\PYG{n}{b}\PYG{p}{)}
\PYG{p}{\PYGZob{}}
\PYG{+w}{  }\PYG{k}{static\PYGZus{}assert}\PYG{p}{(}\PYG{o}{!}\PYG{n}{std}\PYG{o}{:}\PYG{o}{:}\PYG{n}{is\PYGZus{}integral}\PYG{o}{\PYGZlt{}}\PYG{n}{T}\PYG{o}{\PYGZgt{}}\PYG{p}{,}
\PYG{+w}{                }\PYG{l+s}{\PYGZdq{}}\PYG{l+s}{A non\PYGZhy{}integral type expected}\PYG{l+s}{\PYGZdq{}}\PYG{p}{)}\PYG{p}{;}
\PYG{+w}{  }\PYG{k}{return}\PYG{+w}{ }\PYG{n}{a}\PYG{+w}{ }\PYG{o}{*}\PYG{+w}{ }\PYG{n}{b}\PYG{p}{;}
\PYG{p}{\PYGZcb{}}
\PYG{k}{auto}\PYG{+w}{ }\PYG{n}{v1}\PYG{+w}{ }\PYG{o}{=}\PYG{+w}{ }\PYG{n}{compute}\PYG{p}{(}\PYG{l+m+mi}{1}\PYG{p}{,}\PYG{+w}{ }\PYG{l+m+mi}{2}\PYG{p}{)}\PYG{p}{;}
\PYG{c+c1}{// 错误：对重载函数的调用有歧义}
\PYG{k}{auto}\PYG{+w}{ }\PYG{n}{v2}\PYG{+w}{ }\PYG{o}{=}\PYG{+w}{ }\PYG{n}{compute}\PYG{p}{(}\PYG{l+m+mf}{1.0}\PYG{p}{,}\PYG{+w}{ }\PYG{l+m+mf}{2.0}\PYG{p}{)}\PYG{p}{;}
\PYG{c+c1}{// 错误：对重载函数的调用有歧义}
\end{Verbatim}
\end{codeblockbg}


%%C040145 p

无论我们怎么调用这个函数，最终都会出错，因为编译器找到了两个可能被调用的重载。这是因为\texttt{static\_\allowbreak{}assert}只有在重载决议完成之后才会被考虑，而在本例中，重载决议已经构建出一个含两个可能候选的集合。


%%C040146 p

这个问题的解决方案是\texttt{std::\allowbreak{}enable\_\allowbreak{}if}加SFINAE。我们通过前面定义的别名模板\texttt{EnableIf}和\texttt{DisabledIf}，在模板参数上使用\texttt{std::\allowbreak{}enable\_\allowbreak{}if}（尽管第二个重载上我们仍然使用那个占位模板参数，以引出两个不同的定义）。下面的例子展示了改写后的两个重载。第一个重载只对整数类型启用，第二个则对整数类型禁用。


%%C040147 code

\begin{codeblockbg}
\begin{Verbatim}[commandchars=\\\{\},breaklines=true,breakanywhere=true,fontsize=\codefont,xleftmargin=17bp]
\PYG{k}{template}\PYG{+w}{ }\PYG{o}{\PYGZlt{}}\PYG{k}{typename}\PYG{+w}{ }\PYG{n+nc}{T}\PYG{p}{,}\PYG{+w}{ }\PYG{k}{typename}\PYG{+w}{ }\PYG{o}{=}\PYG{+w}{ }\PYG{n}{EnableIf}\PYG{o}{\PYGZlt{}}\PYG{n}{std}\PYG{o}{:}\PYG{o}{:}\PYG{n}{is\PYGZus{}integral}\PYG{o}{\PYGZlt{}}\PYG{n}{T}\PYG{o}{\PYGZgt{}}\PYG{o}{\PYGZgt{}}\PYG{o}{\PYGZgt{}}
\PYG{k}{auto}\PYG{+w}{ }\PYG{n}{compute}\PYG{p}{(}\PYG{n}{T}\PYG{+w}{ }\PYG{k}{const}\PYG{+w}{ }\PYG{n}{a}\PYG{p}{,}\PYG{+w}{ }\PYG{n}{T}\PYG{+w}{ }\PYG{k}{const}\PYG{+w}{ }\PYG{n}{b}\PYG{p}{)}
\PYG{p}{\PYGZob{}}
\PYG{+w}{  }\PYG{k}{return}\PYG{+w}{ }\PYG{n}{a}\PYG{+w}{ }\PYG{o}{*}\PYG{+w}{ }\PYG{n}{b}\PYG{p}{;}
\PYG{p}{\PYGZcb{}}
\PYG{k}{template}\PYG{+w}{ }\PYG{o}{\PYGZlt{}}\PYG{k}{typename}\PYG{+w}{ }\PYG{n+nc}{T}\PYG{p}{,}\PYG{+w}{ }\PYG{k}{typename}\PYG{+w}{ }\PYG{o}{=}\PYG{+w}{ }\PYG{n}{DisableIf}\PYG{o}{\PYGZlt{}}\PYG{n}{std}\PYG{o}{:}\PYG{o}{:}\PYG{n}{is\PYGZus{}integral}\PYG{o}{\PYGZlt{}}\PYG{n}{T}\PYG{o}{\PYGZgt{}}\PYG{o}{\PYGZgt{}}\PYG{p}{,}
\PYG{+w}{          }\PYG{k}{typename}\PYG{+w}{ }\PYG{o}{=}\PYG{+w}{ }\PYG{k+kt}{void}\PYG{o}{\PYGZgt{}}
\PYG{k}{auto}\PYG{+w}{ }\PYG{n}{compute}\PYG{p}{(}\PYG{n}{T}\PYG{+w}{ }\PYG{k}{const}\PYG{+w}{ }\PYG{n}{a}\PYG{p}{,}\PYG{+w}{ }\PYG{n}{T}\PYG{+w}{ }\PYG{k}{const}\PYG{+w}{ }\PYG{n}{b}\PYG{p}{)}
\PYG{p}{\PYGZob{}}
\PYG{+w}{  }\PYG{k}{return}\PYG{+w}{ }\PYG{n}{a}\PYG{+w}{ }\PYG{o}{+}\PYG{+w}{ }\PYG{n}{b}\PYG{p}{;}
\PYG{p}{\PYGZcb{}}
\PYG{k}{auto}\PYG{+w}{ }\PYG{n}{v1}\PYG{+w}{ }\PYG{o}{=}\PYG{+w}{ }\PYG{n}{compute}\PYG{p}{(}\PYG{l+m+mi}{1}\PYG{p}{,}\PYG{+w}{ }\PYG{l+m+mi}{2}\PYG{p}{)}\PYG{p}{;}\PYG{+w}{     }\PYG{c+c1}{// OK；v1 = 2}
\PYG{k}{auto}\PYG{+w}{ }\PYG{n}{v2}\PYG{+w}{ }\PYG{o}{=}\PYG{+w}{ }\PYG{n}{compute}\PYG{p}{(}\PYG{l+m+mf}{1.0}\PYG{p}{,}\PYG{+w}{ }\PYG{l+m+mf}{2.0}\PYG{p}{)}\PYG{p}{;}\PYG{+w}{ }\PYG{c+c1}{// OK；v2 = 3.0}
\end{Verbatim}
\end{codeblockbg}


%%C040148 p

有了SFINAE发挥作用，当编译器为\texttt{compute(\allowbreak{}1,}


%%C040149 li

\begin{itemize}
\item 或\texttt{compute(\allowbreak{}1.\allowbreak{}0, 2.\allowbreak{}0)}构建重载集时；它会直接丢弃那个产生替换

\end{itemize}


%%C040150 p

失败的重载并继续，最终每种情况下重载集都只包含一个候选。


%%C040151 h3

\subsection{延伸阅读}


%%C040152 li

\begin{itemize}
\item \textit{用static\_assert做编译期断言检查}


%%C040153 li

\item \texttt{Chapter 1}，\textit{现代C++} \textit{语言核心特性}中的\textit{创建类型别名与别名模板}案例

\end{itemize}


%%C040154 h2

\section{用constexpr if在编译期选择分支}


%%C040155 p

在前面的案例中，我们看到了如何用\texttt{static\_\allowbreak{}assert}和\texttt{std::\allowbreak{}enable\_\allowbreak{}if}对类型和函数施加限制，以及这两者有何不同。当我们使用SFINAE和\texttt{std::\allowbreak{}enable\_\allowbreak{}if}来定义函数重载，或者编写可变参数函数模板时，模板元编程可能变得复杂而繁杂。C++17的一项新特性正是为了简化这类代码而设计的；它称为\textit{constexpr if}，定义了一种\texttt{if}语句，其条件在编译期求值，结果是由编译器把某个分支的语句体选入翻译单元。constexpr if的典型用途是简化可变参数模板和基于\texttt{std::\allowbreak{}enable\_\allowbreak{}if}的代码。


%%C040156 h3

\subsection{准备工作}


%%C040157 p

本案例将引用并简化前面案例中编写的代码。在继续本案例之前，你应当花点时间回过头看看我们在那些案例中写过的代码，它们如下：


%%C040158 li

\begin{itemize}
\item \textit{用enable\_if条件编译类与函数}案例中针对整数类型和非整数类型的\texttt{compute()}重载。


%%C040159 li

\item \texttt{Chapter 2}，\textit{数字与字符串}中的\textit{创建raw} \textit{用户定义字面量}案例里的用户定义8位、16位和32位二进制字面量。

\end{itemize}


%%C040160 p

这些实现存在几个问题：


%%C040161 li

\begin{itemize}
\item 它们难以阅读。例如，大量注意力都放在模板声明上，而函数体其实非常简单。不过最大的问题在于，它需要开发者付出更多注意力，因为其中充斥着复杂的声明，例如\texttt{typename =} \myttbr{std::\allowbreak{}enable\_\allowbreak{}if<\allowbreak{}std::\allowbreak{}is\_\allowbreak{}integral<\allowbreak{}T>::\allowbreak{}value, T>::\allowbreak{}type}。


%%C040162 li

\item 代码量太大。第一个例子的最终目的是要一个能对不同类型表现出不同行为的泛型函数，可我们却不得不为该函数写两个重载；而且为了区分二者，还得使用一个额外的、未使用的模板参数。第二个例子的目的是用字符\texttt{‘0’}和\texttt{‘1’}构造一个整数值，可我们却不得不写一个类模板和三个特化才能做到。


%%C040163 li

\item 它要求高级的模板元编程技巧，而做这么简单的事情本不该需要这些。

\end{itemize}


%%C040164 p

constexpr if的语法与普通\texttt{if}语句非常相似，只是要求在条件前面加上\texttt{constexpr}关键字。一般形式如下：


%%C040165 code

\begin{codeblockbg}
\begin{Verbatim}[commandchars=\\\{\},breaklines=true,breakanywhere=true,fontsize=\codefont,xleftmargin=17bp]
\PYG{k}{if}\PYG{+w}{ }\PYG{k}{constexpr}\PYG{+w}{ }\PYG{p}{(}\PYG{n}{init}\PYG{o}{\PYGZhy{}}\PYG{n}{statement}\PYG{+w}{ }\PYG{n}{condition}\PYG{p}{)}\PYG{+w}{ }\PYG{n}{statement}\PYG{o}{\PYGZhy{}}\PYG{n+nb}{true}
\PYG{k}{else}\PYG{+w}{ }\PYG{n}{statement}\PYG{o}{\PYGZhy{}}\PYG{n+nb}{false}
\end{Verbatim}
\end{codeblockbg}


%%C040166 h3

\subsection{具体做法}


%%C040167 p

使用constexpr if语句可以做下面这些事：


%%C040168 li

\begin{itemize}
\item 避免使用\texttt{std::\allowbreak{}enable\_\allowbreak{}if}、避免依赖SFINAE来对函数模板类型施加限制，并有条件地编译代码：


%%C040169 code

\begin{codeblockbg}
\begin{Verbatim}[commandchars=\\\{\},breaklines=true,breakanywhere=true,fontsize=\codefont,xleftmargin=17bp]
\PYG{k}{template}\PYG{+w}{ }\PYG{o}{\PYGZlt{}}\PYG{k}{typename}\PYG{+w}{ }\PYG{n+nc}{T}\PYG{o}{\PYGZgt{}}
\PYG{k}{auto}\PYG{+w}{ }\PYG{n}{compute}\PYG{p}{(}\PYG{n}{T}\PYG{+w}{ }\PYG{k}{const}\PYG{+w}{ }\PYG{n}{a}\PYG{p}{,}\PYG{+w}{ }\PYG{n}{T}\PYG{+w}{ }\PYG{k}{const}\PYG{+w}{ }\PYG{n}{b}\PYG{p}{)}
\PYG{p}{\PYGZob{}}
\PYG{+w}{  }\PYG{k}{if}\PYG{+w}{ }\PYG{k}{constexpr}\PYG{+w}{ }\PYG{p}{(}\PYG{n}{std}\PYG{o}{:}\PYG{o}{:}\PYG{n}{is\PYGZus{}integral}\PYG{o}{\PYGZlt{}}\PYG{n}{T}\PYG{o}{\PYGZgt{}}\PYG{o}{:}\PYG{o}{:}\PYG{n}{value}\PYG{p}{)}
\PYG{+w}{    }\PYG{k}{return}\PYG{+w}{ }\PYG{n}{a}\PYG{+w}{ }\PYG{o}{*}\PYG{+w}{ }\PYG{n}{b}\PYG{p}{;}
\PYG{+w}{  }\PYG{k}{else}
\PYG{+w}{    }\PYG{k}{return}\PYG{+w}{ }\PYG{n}{a}\PYG{+w}{ }\PYG{o}{+}\PYG{+w}{ }\PYG{n}{b}\PYG{p}{;}
\PYG{p}{\PYGZcb{}}
\end{Verbatim}
\end{codeblockbg}


%%C040170 li

\item 简化可变参数模板的编写，并实现元编程式的编译期递归：


%%C040171 code

\begin{codeblockbg}
\begin{Verbatim}[commandchars=\\\{\},breaklines=true,breakanywhere=true,fontsize=\codefont,xleftmargin=17bp]
\PYG{k}{namespace}\PYG{+w}{ }\PYG{n+nn}{binary}
\PYG{p}{\PYGZob{}}
\PYG{+w}{  }\PYG{k}{using}\PYG{+w}{ }\PYG{n}{byte8}\PYG{+w}{ }\PYG{o}{=}\PYG{+w}{ }\PYG{k+kt}{unsigned}\PYG{+w}{ }\PYG{k+kt}{char}\PYG{p}{;}
\PYG{k}{namespace}\PYG{+w}{ }\PYG{n+nn}{binary\PYGZus{}literals}
\PYG{p}{\PYGZob{}}
\PYG{+w}{  }\PYG{k}{namespace}\PYG{+w}{ }\PYG{n+nn}{binary\PYGZus{}literals\PYGZus{}internals}
\PYG{+w}{  }\PYG{p}{\PYGZob{}}
\PYG{+w}{    }\PYG{k}{template}\PYG{+w}{ }\PYG{o}{\PYGZlt{}}\PYG{k}{typename}\PYG{+w}{ }\PYG{n+nc}{CharT}\PYG{p}{,}\PYG{+w}{ }\PYG{k+kt}{char}\PYG{+w}{ }\PYG{n}{d}\PYG{p}{,}\PYG{+w}{ }\PYG{k+kt}{char}\PYG{p}{.}\PYG{p}{.}\PYG{p}{.}\PYG{+w}{ }\PYG{n}{bits}\PYG{o}{\PYGZgt{}}
\PYG{+w}{    }\PYG{k}{constexpr}\PYG{+w}{ }\PYG{n}{CharT}\PYG{+w}{ }\PYG{n}{binary\PYGZus{}eval}\PYG{p}{(}\PYG{p}{)}
\PYG{+w}{    }\PYG{p}{\PYGZob{}}
\PYG{+w}{      }\PYG{k}{if}\PYG{+w}{ }\PYG{k}{constexpr}\PYG{p}{(}\PYG{k}{sizeof}\PYG{p}{.}\PYG{p}{.}\PYG{p}{.}\PYG{p}{(}\PYG{n}{bits}\PYG{p}{)}\PYG{+w}{ }\PYG{o}{=}\PYG{o}{=}\PYG{+w}{ }\PYG{l+m+mi}{0}\PYG{p}{)}
\PYG{+w}{        }\PYG{k}{return}\PYG{+w}{ }\PYG{k}{static\PYGZus{}cast}\PYG{o}{\PYGZlt{}}\PYG{n}{CharT}\PYG{o}{\PYGZgt{}}\PYG{p}{(}\PYG{n}{d}\PYG{o}{\PYGZhy{}}\PYG{l+s+sc}{\PYGZsq{}}\PYG{l+s+sc}{0}\PYG{l+s+sc}{\PYGZsq{}}\PYG{p}{)}\PYG{p}{;}
\PYG{+w}{      }\PYG{k}{else}\PYG{+w}{ }\PYG{k}{if}\PYG{+w}{ }\PYG{k}{constexpr}\PYG{p}{(}\PYG{n}{d}\PYG{+w}{ }\PYG{o}{=}\PYG{o}{=}\PYG{+w}{ }\PYG{l+s+sc}{\PYGZsq{}}\PYG{l+s+sc}{0}\PYG{l+s+sc}{\PYGZsq{}}\PYG{p}{)}
\PYG{+w}{        }\PYG{k}{return}\PYG{+w}{ }\PYG{n}{binary\PYGZus{}eval}\PYG{o}{\PYGZlt{}}\PYG{n}{CharT}\PYG{p}{,}\PYG{+w}{ }\PYG{n}{bits}\PYG{p}{.}\PYG{p}{.}\PYG{p}{.}\PYG{o}{\PYGZgt{}}\PYG{p}{(}\PYG{p}{)}\PYG{p}{;}
\PYG{+w}{      }\PYG{k}{else}\PYG{+w}{ }\PYG{k}{if}\PYG{+w}{ }\PYG{k}{constexpr}\PYG{p}{(}\PYG{n}{d}\PYG{+w}{ }\PYG{o}{=}\PYG{o}{=}\PYG{+w}{ }\PYG{l+s+sc}{\PYGZsq{}}\PYG{l+s+sc}{1}\PYG{l+s+sc}{\PYGZsq{}}\PYG{p}{)}
\PYG{+w}{        }\PYG{k}{return}\PYG{+w}{ }\PYG{k}{static\PYGZus{}cast}\PYG{o}{\PYGZlt{}}\PYG{n}{CharT}\PYG{o}{\PYGZgt{}}\PYG{p}{(}
\PYG{+w}{          }\PYG{p}{(}\PYG{l+m+mi}{1}\PYG{+w}{ }\PYG{o}{\PYGZlt{}}\PYG{o}{\PYGZlt{}}\PYG{+w}{ }\PYG{k}{sizeof}\PYG{p}{.}\PYG{p}{.}\PYG{p}{.}\PYG{p}{(}\PYG{n}{bits}\PYG{p}{)}\PYG{p}{)}\PYG{+w}{ }\PYG{o}{|}
\PYG{+w}{          }\PYG{n}{binary\PYGZus{}eval}\PYG{o}{\PYGZlt{}}\PYG{n}{CharT}\PYG{p}{,}\PYG{+w}{ }\PYG{n}{bits}\PYG{p}{.}\PYG{p}{.}\PYG{p}{.}\PYG{o}{\PYGZgt{}}\PYG{p}{(}\PYG{p}{)}\PYG{p}{)}\PYG{p}{;}
\PYG{+w}{    }\PYG{p}{\PYGZcb{}}
\PYG{+w}{  }\PYG{p}{\PYGZcb{}}
\PYG{k}{template}\PYG{o}{\PYGZlt{}}\PYG{k+kt}{char}\PYG{p}{.}\PYG{p}{.}\PYG{p}{.}\PYG{+w}{ }\PYG{n}{bits}\PYG{o}{\PYGZgt{}}
\PYG{k}{constexpr}\PYG{+w}{ }\PYG{n}{byte8}\PYG{+w}{ }\PYG{k}{operator}\PYG{l+s}{\PYGZdq{}}\PYG{l+s}{\PYGZdq{}}\PYG{n}{\PYGZus{}b8}\PYG{p}{(}\PYG{p}{)}
\PYG{p}{\PYGZob{}}
\PYG{+w}{  }\PYG{k}{static\PYGZus{}assert}\PYG{p}{(}
\PYG{+w}{    }\PYG{k}{sizeof}\PYG{p}{.}\PYG{p}{.}\PYG{p}{.}\PYG{p}{(}\PYG{n}{bits}\PYG{p}{)}\PYG{+w}{ }\PYG{o}{\PYGZlt{}}\PYG{o}{=}\PYG{+w}{ }\PYG{l+m+mi}{8}\PYG{p}{,}
\PYG{+w}{    }\PYG{l+s}{\PYGZdq{}}\PYG{l+s}{binary literal b8 must be up to 8 digits long}\PYG{l+s}{\PYGZdq{}}\PYG{p}{)}\PYG{p}{;}
\PYG{+w}{      }\PYG{k}{return}\PYG{+w}{ }\PYG{n}{binary\PYGZus{}literals\PYGZus{}internals}\PYG{o}{:}\PYG{o}{:}
\PYG{+w}{                  }\PYG{n}{binary\PYGZus{}eval}\PYG{o}{\PYGZlt{}}\PYG{n}{byte8}\PYG{p}{,}\PYG{+w}{ }\PYG{n}{bits}\PYG{p}{.}\PYG{p}{.}\PYG{p}{.}\PYG{o}{\PYGZgt{}}\PYG{p}{(}\PYG{p}{)}\PYG{p}{;}
\PYG{+w}{    }\PYG{p}{\PYGZcb{}}
\PYG{+w}{  }\PYG{p}{\PYGZcb{}}
\PYG{p}{\PYGZcb{}}
\end{Verbatim}
\end{codeblockbg}

\end{itemize}


%%C040172 h3

\subsection{工作原理}


%%C040173 p

constexpr if的工作方式相对简单：\texttt{if}语句中的条件必须是一个编译期表达式，它求值的结果可以转换为布尔值。如果条件为\texttt{true}，就选择\texttt{if}语句的语句体，也就是说它会进入翻译单元参与编译。如果条件为\texttt{false}，则求值\texttt{else}分支（如果定义了的话）。被丢弃的constexpr if分支中的return语句不参与函数返回类型推导。


%%C040174 p

在\textit{具体做法}一节的第一个例子中，\texttt{compute()}函数模板的签名很干净，函数体也非常简单；如果替换模板参数的类型是整数类型，编译器就会选择第一个分支（即\texttt{return a *} \texttt{b;}）用于代码生成，并丢弃\texttt{else}分支。对于非整数类型，由于条件求值为\texttt{false}，编译器会选择\texttt{else}分支（即\texttt{return a +} \texttt{b;}）用于代码生成，并丢弃其余部分。

%%C040175 p

在\textit{具体做法}一节的第二个例子中，内部的辅助函数\texttt{binary\_\allowbreak{}eval()}是一个不带任何参数的可变参数模板函数；它只有模板参数。该函数先对第一个参数求值，再以递归方式处理其余参数（但要记住，这并不是运行期递归）。当只剩下一个字符、且剩余参数包的大小为0时，我们返回该字符所表示的十进制值（\texttt{‘0’}对应0，\texttt{‘1’}对应1）。如果当前第一个元素是\texttt{‘0’}，我们返回对剩余参数包求值所得的结果，这涉及一次递归调用。如果当前第一个元素是\texttt{‘1’}，我们返回的值要么是把1向左移动若干位（位数由剩余参数包的大小决定）所得的结果，要么是对剩余参数包求值所得的结果，后者同样涉及一次递归调用。


%%C040176 h2

\section{用属性向编译器提供元数据}


%%C040177 p

在涉及类型或数据的反射与内省特性，以及定义语言扩展的标准机制方面，C++一直非常欠缺。正因如此，各家编译器为此定义了自己的专用扩展，例如VC++的\texttt{\_\_\allowbreak{}declspec()}说明符或GCC的\myttbr{\_\_\allowbreak{}attribute\_\_((...))}。不过，C++11引入了属性的概念，使编译器能够以标准方式实现扩展，甚至实现内嵌的领域特定语言。新的C++标准定义了几个所有编译器都应实现的属性，这就是本案例的主题。


%%C040178 h3

\subsection{具体做法}


%%C040179 p

使用标准属性向编译器提示各种设计目标：


%%C040180 li

\begin{itemize}
\item 为了确保函数的返回值不能被忽略，用\texttt{[[nodiscard]]}属性声明该函数：


%%C040181 code

\begin{codeblockbg}
\begin{Verbatim}[commandchars=\\\{\},breaklines=true,breakanywhere=true,fontsize=\codefont,xleftmargin=17bp]
\PYG{p}{[[}\PYG{n+na}{nodiscard}\PYG{p}{]]}\PYG{+w}{ }\PYG{k+kt}{int}\PYG{+w}{ }\PYG{n}{get\PYGZus{}value1}\PYG{p}{(}\PYG{p}{)}
\PYG{p}{\PYGZob{}}
\PYG{+w}{  }\PYG{k}{return}\PYG{+w}{ }\PYG{l+m+mi}{42}\PYG{p}{;}
\PYG{p}{\PYGZcb{}}
\PYG{n}{get\PYGZus{}value1}\PYG{p}{(}\PYG{p}{)}\PYG{p}{;}
\PYG{c+c1}{// 警告：忽略了函数的返回值}
\PYG{c+c1}{//          该函数用 \PYGZsq{}nodiscard\PYGZsq{} 属性声明 get\PYGZus{}value1();}
\end{Verbatim}
\end{codeblockbg}


%%C040182 li

\item 或者，你也可以用\texttt{[[nodiscard]]}属性声明用作函数返回类型的枚举和类；在这种情况下，任何返回这类类型的函数的返回值都不能被忽略：


%%C040183 code

\begin{codeblockbg}
\begin{Verbatim}[commandchars=\\\{\},breaklines=true,breakanywhere=true,fontsize=\codefont,xleftmargin=17bp]
\PYG{k}{enum}\PYG{+w}{ }\PYG{k}{class}\PYG{p}{[[}\PYG{n+na}{nodiscard}\PYG{p}{]]}\PYG{+w}{ }\PYG{n+nc}{ReturnCodes}\PYG{p}{\PYGZob{}}\PYG{+w}{ }\PYG{n}{OK}\PYG{p}{,}\PYG{+w}{ }\PYG{n}{NoData}\PYG{p}{,}\PYG{+w}{ }\PYG{n}{Error}\PYG{+w}{ }\PYG{p}{\PYGZcb{}}\PYG{p}{;}
\PYG{n}{ReturnCodes}\PYG{+w}{ }\PYG{n+nf}{get\PYGZus{}value2}\PYG{p}{(}\PYG{p}{)}
\PYG{p}{\PYGZob{}}
\PYG{+w}{  }\PYG{k}{return}\PYG{+w}{ }\PYG{n}{ReturnCodes}\PYG{o}{:}\PYG{o}{:}\PYG{n}{OK}\PYG{p}{;}
\PYG{p}{\PYGZcb{}}
\PYG{k}{struct}\PYG{p}{[[}\PYG{n+na}{nodiscard}\PYG{p}{]]}\PYG{+w}{ }\PYG{n}{Item}\PYG{p}{\PYGZob{}}\PYG{p}{\PYGZcb{}}\PYG{p}{;}
\PYG{n}{Item}\PYG{+w}{ }\PYG{n+nf}{get\PYGZus{}value3}\PYG{p}{(}\PYG{p}{)}
\PYG{p}{\PYGZob{}}
\PYG{+w}{  }\PYG{k}{return}\PYG{+w}{ }\PYG{n}{Item}\PYG{p}{\PYGZob{}}\PYG{p}{\PYGZcb{}}\PYG{p}{;}
\PYG{p}{\PYGZcb{}}
\PYG{c+c1}{// 警告：忽略了函数的返回值}
\PYG{c+c1}{//          该函数用 \PYGZsq{}nodiscard\PYGZsq{} 属性声明}
\PYG{n}{get\PYGZus{}value2}\PYG{p}{(}\PYG{p}{)}\PYG{p}{;}
\PYG{n}{get\PYGZus{}value3}\PYG{p}{(}\PYG{p}{)}\PYG{p}{;}
\end{Verbatim}
\end{codeblockbg}


%%C040184 li

\item 为了确保使用被视为已弃用的函数或类型时编译器会给出警告，用\texttt{[[deprecated]]}属性声明它们：


%%C040185 code

\begin{codeblockbg}
\begin{Verbatim}[commandchars=\\\{\},breaklines=true,breakanywhere=true,fontsize=\codefont,xleftmargin=17bp]
\PYG{p}{[[}\PYG{n+na}{deprecated}\PYG{p}{(}\PYG{l+s}{\PYGZdq{}}\PYG{l+s}{Use func2()}\PYG{l+s}{\PYGZdq{}}\PYG{p}{)}\PYG{p}{]]}\PYG{+w}{ }\PYG{k+kt}{void}\PYG{+w}{ }\PYG{n}{func}\PYG{p}{(}\PYG{p}{)}
\PYG{p}{\PYGZob{}}
\PYG{p}{\PYGZcb{}}
\PYG{c+c1}{// 警告：\PYGZsq{}func\PYGZsq{} 已弃用：请使用 func2()}
\PYG{n}{func}\PYG{p}{(}\PYG{p}{)}\PYG{p}{;}
\PYG{k}{class}\PYG{+w}{ }\PYG{p}{[[}\PYG{n+na}{deprecated}\PYG{p}{]]}\PYG{+w}{ }\PYG{n+nc}{foo}
\PYG{p}{\PYGZob{}}
\PYG{p}{\PYGZcb{}}\PYG{p}{;}
\PYG{c+c1}{// 警告：\PYGZsq{}foo\PYGZsq{} 已弃用}
\PYG{n}{foo}\PYG{+w}{ }\PYG{n}{f}\PYG{p}{;}
\end{Verbatim}
\end{codeblockbg}


%%C040186 li

\item 为了确保编译器不为未使用的变量发出警告，使用\texttt{[[maybe\_\allowbreak{}unused]]}属性：


%%C040187 code

\begin{codeblockbg}
\begin{Verbatim}[commandchars=\\\{\},breaklines=true,breakanywhere=true,fontsize=\codefont,xleftmargin=17bp]
\PYG{k+kt}{double}\PYG{+w}{ }\PYG{n+nf}{run}\PYG{p}{(}\PYG{p}{[[}\PYG{n+na}{maybe\PYGZus{}unused}\PYG{p}{]]}\PYG{+w}{ }\PYG{k+kt}{int}\PYG{+w}{ }\PYG{n}{a}\PYG{p}{,}\PYG{+w}{ }\PYG{k+kt}{double}\PYG{+w}{ }\PYG{n}{b}\PYG{p}{)}
\PYG{p}{\PYGZob{}}
\PYG{+w}{  }\PYG{k}{return}\PYG{+w}{ }\PYG{l+m+mi}{2}\PYG{+w}{ }\PYG{o}{*}\PYG{+w}{ }\PYG{n}{b}\PYG{p}{;}
\PYG{p}{\PYGZcb{}}
\PYG{p}{[[}\PYG{n+na}{maybe\PYGZus{}unused}\PYG{p}{]]}\PYG{+w}{ }\PYG{k}{auto}\PYG{+w}{ }\PYG{n}{i}\PYG{+w}{ }\PYG{o}{=}\PYG{+w}{ }\PYG{n}{get\PYGZus{}value1}\PYG{p}{(}\PYG{p}{)}\PYG{p}{;}
\end{Verbatim}
\end{codeblockbg}


%%C040188 li

\item 为了确保\texttt{switch}语句中有意为之的贯穿（fall-through）case标签不被编译器警告，使用\texttt{[[fallthrough]]}属性：


%%C040189 code

\begin{codeblockbg}
\begin{Verbatim}[commandchars=\\\{\},breaklines=true,breakanywhere=true,fontsize=\codefont,xleftmargin=17bp]
\PYG{k+kt}{void}\PYG{+w}{ }\PYG{n+nf}{option1}\PYG{p}{(}\PYG{p}{)}\PYG{+w}{ }\PYG{p}{\PYGZob{}}\PYG{p}{\PYGZcb{}}
\PYG{k+kt}{void}\PYG{+w}{ }\PYG{n+nf}{option2}\PYG{p}{(}\PYG{p}{)}\PYG{+w}{ }\PYG{p}{\PYGZob{}}\PYG{p}{\PYGZcb{}}
\PYG{k+kt}{int}\PYG{+w}{ }\PYG{n}{alternative}\PYG{+w}{ }\PYG{o}{=}\PYG{+w}{ }\PYG{n}{get\PYGZus{}value1}\PYG{p}{(}\PYG{p}{)}\PYG{p}{;}
\PYG{k}{switch}\PYG{+w}{ }\PYG{p}{(}\PYG{n}{alternative}\PYG{p}{)}
\PYG{p}{\PYGZob{}}
\PYG{+w}{  }\PYG{k}{case}\PYG{+w}{ }\PYG{l+m+mi}{1}\PYG{p}{:}
\PYG{+w}{    }\PYG{n}{option1}\PYG{p}{(}\PYG{p}{)}\PYG{p}{;}
\PYG{+w}{    }\PYG{p}{[[}\PYG{n+na}{fallthrough}\PYG{p}{]]}\PYG{p}{;}\PYG{+w}{ }\PYG{c+c1}{// 这是有意为之}
\PYG{+w}{  }\PYG{k}{case}\PYG{+w}{ }\PYG{l+m+mi}{2}\PYG{p}{:}
\PYG{+w}{   }\PYG{n}{option2}\PYG{p}{(}\PYG{p}{)}\PYG{p}{;}
\PYG{p}{\PYGZcb{}}
\end{Verbatim}
\end{codeblockbg}

\end{itemize}


%%C040190 h3

\subsection{工作原理}


%%C040191 p

属性是C++的一项非常灵活的特性；它们几乎可以用在任何地方，但具体用法是针对每个属性单独定义的。它们可以用于类型、函数、变量、名字、代码块乃至整个翻译单元。


%%C040192 p

属性写在双方括号之间（例如\texttt{[[attr1]]}），并且在一个声明中可以指定多个属性（例如\texttt{[[attr1, attr2,} \texttt{attr3]]}）。


%%C040193 p

属性可以带参数，例如\texttt{[[mode(\allowbreak{}greedy)\allowbreak{}]]}；也可以使用完整限定名，例如\texttt{[[sys::\allowbreak{}hidden]]}或\myttbr{[[using sys: visibility(\allowbreak{}hidden)\allowbreak{},} \texttt{debug]]}。


%%C040194 p

属性既可以出现在所应用实体的名字之前，也可以出现在其后，还可以两处都出现，此时它们会被合并。下面给出几个说明这一点的例子：


%%C040195 code

\begin{codeblockbg}
\begin{Verbatim}[commandchars=\\\{\},breaklines=true,breakanywhere=true,fontsize=\codefont,xleftmargin=17bp]
\PYG{c+c1}{// attr1 作用于 a，attr2 作用于 b}
\PYG{k+kt}{int}\PYG{+w}{ }\PYG{n}{a}\PYG{+w}{ }\PYG{p}{[[}\PYG{n+na}{attr1}\PYG{p}{]]}\PYG{p}{,}\PYG{+w}{ }\PYG{n}{b}\PYG{+w}{ }\PYG{p}{[[}\PYG{n+na}{attr2}\PYG{p}{]]}\PYG{p}{;}
\PYG{c+c1}{// attr1 作用于 a 和 b}
\PYG{k+kt}{int}\PYG{+w}{ }\PYG{p}{[[}\PYG{n+na}{attr1}\PYG{p}{]]}\PYG{+w}{ }\PYG{n}{a}\PYG{p}{,}\PYG{+w}{ }\PYG{n}{b}\PYG{p}{;}
\PYG{c+c1}{// attr1 作用于 a 和 b，attr2 作用于 a}
\PYG{k+kt}{int}\PYG{+w}{ }\PYG{p}{[[}\PYG{n+na}{attr1}\PYG{p}{]]}\PYG{+w}{ }\PYG{n}{a}\PYG{+w}{ }\PYG{p}{[[}\PYG{n+na}{attr2}\PYG{p}{]]}\PYG{p}{,}\PYG{+w}{ }\PYG{n}{b}\PYG{p}{;}
\end{Verbatim}
\end{codeblockbg}


%%C040196 p

属性不能出现在命名空间声明中，但可以作为单行声明出现在命名空间中的任何位置。此时，它究竟作用于后面的声明、作用于命名空间，还是作用于翻译单元，取决于具体的属性：


%%C040197 code

\begin{codeblockbg}
\begin{Verbatim}[commandchars=\\\{\},breaklines=true,breakanywhere=true,fontsize=\codefont,xleftmargin=17bp]
\PYG{k}{namespace}\PYG{+w}{ }\PYG{n+nn}{test}
\PYG{p}{\PYGZob{}}
\PYG{+w}{  }\PYG{p}{[[}\PYG{n+na}{debug}\PYG{p}{]]}\PYG{p}{;}
\PYG{p}{\PYGZcb{}}
\end{Verbatim}
\end{codeblockbg}


%%C040198 p

在现代C++编程的书籍和教程中，属性常常被忽略或只是一笔带过，原因大概是开发者实际上无法自己编写属性，因为这一语言特性是面向编译器实现者的。不过，对某些编译器来说，定义用户提供的属性或许是可行的；GCC就是这样的编译器，它支持通过插件为编译器添加额外功能，这些插件也可以用来定义新的属性。


%%C040199 p

不过，标准确实定义了几个所有编译器都必须实现的属性，使用它们有助于你写出更好的代码。我们在上一节的例子中已经见过其中一些。这些属性是在不同版本的标准中定义的：


%%C040200 li

\begin{itemize}
\item 在C++11中：


%%C040201 li

\begin{itemize}
\item \texttt{[[noreturn]]}属性表示函数不会返回。


%%C040202 li

\item \myttbr{[[carries\_\allowbreak{}dependency]]}属性表示release-consume的\texttt{std::\allowbreak{}memory\_\allowbreak{}order}中的依赖链会从函数传入并传出，这使编译器可以跳过不必要的内存栅栏指令。


%%C040203 li

\item 在C++14中：


%%C040204 li

\item \texttt{[[deprecated]]}和\myttbr{[[deprecated(\allowbreak{}“reason”)\allowbreak{}]]}属性表示用这些属性声明的实体被视为已弃用，不应再使用。这些属性可用于类、非静态数据成员、typedef、函数、枚举和模板特化。原因字符串是可选参数。


%%C040205 li

\item 在C++17中：


%%C040206 li

\item \texttt{[[fallthrough]]}属性表示\texttt{switch}语句中标签之间的贯穿是有意为之的。该属性必须单独占一行，紧接在某个case标签之前出现。


%%C040207 li

\item \texttt{[[nodiscard]]}属性表示函数的返回值不能被忽略。


%%C040208 li

\item \texttt{[[maybe\_\allowbreak{}unused]]}属性表示某个实体可能未被使用，但编译器不应对此发出警告。该属性可应用于变量、类、非静态数据成员、枚举、枚举项和typedef。


\end{itemize}


\end{itemize}
